% concatenated from main.tex, tables/tab01_base_case.tex, tables/tab01_commodity_outcomes.tex, tables/tab02_voss.tex, tables/tab03_vmc.tex, tables/tab13_stress_test.tex, tables/tab04_vdu.tex, tables/tab12_vep_sweep.tex, tables/tab05_frontier.tex, tables/tab14_alpha_sweep.tex, tables/tab15_scenario_scaling.tex, tables/tab10_benders_convergence.tex, appendix.tex, tables/tab01_scenario_library.tex, tables/tab01_scenario_outcomes.tex, tables/tab02_design_comparison.tex, tables/tab03_probability_comparison.tex, tables/tab04_design_comparison.tex
\documentclass[final,3p,authoryear]{elsarticle}


\journal{Transportation Research Part E}
% \usepackage{lineno}
% \modulolinenumbers[1]


\usepackage{amsmath, amssymb, amsthm, mathtools}
\usepackage{booktabs}
\usepackage{array}
\usepackage{multirow}
\usepackage{float}
\usepackage{graphicx}
\usepackage{xcolor}
\usepackage{xurl}
\usepackage{enumitem}
\usepackage{setspace}
\usepackage{caption}
\usepackage{subcaption}
\usepackage{longtable}
\usepackage{tabularx}
\newcolumntype{C}{>{\centering\arraybackslash}X}
\usepackage{threeparttable}
\usepackage{algorithm}
\usepackage{rotating}
\usepackage{algpseudocode}
\usepackage{tcolorbox}
\usepackage{adjustbox}
\usepackage{makecell} 
\usepackage{xltabular} 
\usepackage{pdflscape}


\usepackage{hyperref}
\hypersetup{
    colorlinks=true,
    linkcolor=blue,
    citecolor=blue,
    urlcolor=blue
}
\usepackage[nameinlink,capitalise,noabbrev]{cleveref}

\onehalfspacing
\setlength{\parindent}{1.5em}
\setlength{\parskip}{0.3em}
\renewcommand{\arraystretch}{1.15}


\newtheorem{theorem}{Theorem}
\newtheorem{proposition}{Proposition}
\newtheorem{lemma}[theorem]{Lemma}
\newtheorem{corollary}[theorem]{Corollary}
\theoremstyle{definition}
\newtheorem{definition}{Definition}
\newtheorem{assumption}{Assumption}
\newtheorem{remark}{Remark}


\newcommand{\E}{\mathbb{E}}
\newcommand{\Prob}{\mathbb{P}}
\newcommand{\R}{\mathbb{R}}
\newcommand{\Z}{\mathbb{Z}}
\DeclareMathOperator{\cvar}{CVaR}
\DeclareMathOperator{\var}{VaR}
\DeclareMathOperator*{\argmin}{arg\,min}
\newcommand{\arc}[2]{(#1,#2)}


\begin{document}

\begin{frontmatter}

\title{Securing the Flow: Maritime Energy Resilience under Correlated and Decision-Dependent Disruptions}

%% Group authors per Elsevier's style
\author[addr1]{Monit Sharma}
% \ead{monitsharma@smu.edu.sg}

\author[addr1]{Hoong Chuin Lau\corref{cor1}}
\ead{hclau@smu.edu.sg}

\address[addr1]{School of Computing and Information Systems,  Singapore Management University, Singapore}
% \address[addr2]{Institute of High Performance Computing, A*STAR, Singapore}

\cortext[cor1]{Corresponding author}

\begin{abstract}
We develop a two-stage stochastic multi-commodity flow model for
designing a resilient maritime energy supply network under correlated
chokepoint disruptions. A national energy-security planner selects
strategic inventory positions and infrastructure activations before
uncertainty is resolved, then routes four commodity classes—crude oil,
LNG, LPG, and fertilizer—through a capacitated network as recourse.
Three features distinguish the model from the existing literature:
disruption scenarios across maritime corridors are \emph{correlated},
reflecting the empirical reality that geographically proximate
chokepoints (Hormuz, Bab~el-Mandeb) fail jointly; scenario probabilities
depend endogenously on first-stage design decisions via an affine
distortion, formalizing \emph{risk exposure through utilization}; and a
mean--CVaR objective protects against tail-risk shortage outcomes.
We establish that the decision-dependent probability model admits an
exact MILP reformulation via McCormick linearization, that CVaR
preserves scenario-wise decomposability, and that the resulting Benders
decomposition with corridor-based group-failure cuts converges in
finitely many iterations. The model is calibrated to Indian maritime
energy imports using EIA, UNCTAD, World Bank, and operational data from
the 2026 Hormuz crisis. On the calibrated 16-node, 28-arc instance, Benders recovers the extensive-form optimum to numerical precision at every scenario size up to $|S|=729$, with iteration count constant at 10--11; the empirical content of the finite-convergence result holds along the scenario dimension On the calibrated network, the value of stochastic modeling (VSS) is 14.8\% of the optimal objective; the value of
decision-dependent probabilities (VEP) ranges from 0.93\% at the
analyst-set exposure sensitivity to 8.18\% at the probability-validity
boundary; the mean--CVaR frontier exhibits a design phase transition at
confidence level $\alpha\approx 0.75$; and the value of modeling
correlation is identically zero across all $\rho$- and
$\gamma$-stress tests, a substantive finding that we interpret as
\emph{structural joint-failure resilience} of the calibrated network:
the diversified route-and-inventory portfolio selected by the
stochastic program absorbs joint-corridor disruption through the same
hedging mechanism that absorbs single-corridor disruption. LPG emerges
as the most exposed commodity under the optimal policy, while crude oil
is fully hedgeable via pipeline bypass and strategic reserves.
\end{abstract}

\begin{keyword}
Maritime Logistics \sep Supply Chain Resilience \sep Stochastic Programming \sep Decision-Dependent Probabilities \sep Chokepoint Disruptions
\end{keyword}

\end{frontmatter}



\section{Introduction}
\label{sec:introduction}
%================================================================

Capacitated network design under uncertainty is a foundational challenge in
operations research, with applications spanning telecommunications, freight
logistics, humanitarian relief, and energy supply chains. When disruptions
can disable network arcs, the planner faces a two-stage decision: invest in
infrastructure, backup capacity, and strategic inventory \emph{before}
uncertainty is resolved, then route commodities through the surviving
network \emph{after} a disruption scenario materializes. 
While the Stochastic Capacitated Multimodal Network Design Problem (S-CMNDP) has been studied extensively, existing models often rely on the simplifying and frequently inaccurate assumptions of independent failures and exogenous disruption probabilities \citep{snyder2005, lim2010}. 

In this paper, we consider the resilience of an import-dependent economy faced with maritime disruptions, and address three interlocking and practically important extensions to S-CMNDP, namely correlated geographic risk, endogenous exposure, and tail-risk sensitivity. We propose to combine them in a single tractable model, and show that their interaction produces qualitatively different design prescriptions.

To capture the strategic reality of maritime energy supply resilience, our model incorporates a dual-risk framework that distinguishes between how risks are structured in the environment and how they are triggered by the planner’s actions:
\begin{enumerate}
    \item Correlated Environmental Risk: We move beyond the assumption of independent arc failures by modeling geographic correlation. This reflects the empirical reality that maritime chokepoints are not isolated; a disruption in the Strait of Hormuz significantly elevates the simultaneous risk at the Bab el-Mandeb due to shared regional geopolitical friction.
    \item Decision-Dependent Behavioral Risk: Conversely, we recognize that risk is not purely exogenous. Using decision-dependent probabilities, the model acknowledges that a planner’s commitment to a specific corridor—through long-term contracts or infrastructure activation—endogenously shifts the nation’s risk exposure. In this framework, the 'act of routing' is itself a risk-modulating decision, as high utilization of vulnerable paths increases the probability mass of disruption scenarios in the planner’s strategic assessment. 
\end{enumerate}

\paragraph{Motivating Problem}
Maritime chokepoints provide a compelling application for this model. The Strait of Hormuz, the Bab el-Mandeb strait, and the Suez Canal collectively channel over 40\% of globally traded crude oil and liquefied natural gas (LNG).
These corridors are geographically proximate and geopolitically linked:
disruption at one raises the likelihood of cascading disruption at
neighboring passages. For an import dependent economy like India, the world's third-largest energy consumer which
imports approximately 4.5~million barrels per day of crude oil from
Gulf producers via Hormuz, alongside 25~Mt/yr of LNG (50\% from
Qatar), 18~Mt/yr of LPG (90\% from the Middle East), and 5~Mt/yr of
fertilizer feedstock, such disruptions create a design challenge. 
A planner must design
its supply network, selecting suppliers, transport modes, bypass
infrastructure, and strategic stockholding, against a correlated
disruption landscape in which the very act of routing through a vulnerable
corridor contributes to the nation's risk exposure. This necessitates a model that captures (i)~multi-commodity flows with commodity-specific
admissibility constraints, (ii)~correlated disruption scenarios calibrated
to real geopolitical data, (iii)~endogenous probability distortion via
affine decision-dependence, and (iv)~tail-risk protection through a mean--CVaR objective.



\paragraph{Contributions}
This paper makes three contributions to the stochastic network design literature.

\begin{enumerate}[leftmargin=2em, itemsep=0.3em]
    \item \textbf{Formulation.} We present a two-stage stochastic model
    for capacitated multi-commodity network design that simultaneously
    incorporates commodity-specific admissibility, correlated arc
    failures across disruption groups, decision-dependent scenario
    probabilities via affine distortion, corridor-dependence caps, and
    a mean--CVaR objective. We prove that the resulting bilinear model
    admits an exact MILP reformulation
    (\Cref{prop:milp-reform}) via McCormick linearization,
    and that the CVaR term preserves scenario-wise decomposability
    (\Cref{prop:cvar-decomp}). We are not aware of a prior prescriptive
    maritime energy network design model that combines all five features
    in a single tractable formulation.

    \item \textbf{Algorithm.} We develop a Benders decomposition with
    correlation-exploiting corridor-based group-failure cuts and prove
    finite convergence (\Cref{prop:benders_finite}). The implementation
    matches the extensive-form objective to machine precision on every
    converged instance and recovers identical first-stage solutions.
    On the calibrated network, scenario count grows from $|S|=9$ to
    $|S|=729$ with iteration count essentially constant at 10--11,
    confirming the empirical content of the convergence result.

    \item \textbf{Computational study and managerial findings.}
    Calibrating the model to the Indian maritime energy import network
using EIA, UNCTAD, World Bank, and 2026 Hormuz-crisis operational data
(protocol in \ref{app:calibration}) yields three insights. First, the
decision-dependent probability channel reshapes the first-stage route
portfolio away from utilization-sensitive corridors, with the magnitude
scaling monotonically in the exposure-sensitivity prior~$\delta$; the
calibrated network further exhibits \emph{structural joint-failure
resilience}, identifying~$\delta$ rather than the joint-probability
structure as the central calibration object. Second, crude oil and
fertilizer are fully protected through pipeline bypass and strategic
reserves, while LPG emerges as the most exposed commodity and the
natural target of incremental resilience investment. Third, a joint
$(\lambda,\alpha)$ sweep of the mean--CVaR objective traces the
planner's price of resilience and resolves the apparent flatness of
the $\lambda$-only frontier as a discrete tail-saturation artifact,
with the Benders implementation remaining numerically exact across the
scenario libraries tested.


\end{enumerate}



\paragraph{Assumptions}
Our formulation assumes that there is a coordinated national-level energy security agency that exercises four distinct authorities: strategic
reserve holding, infrastructure enablement (pipeline sanctioning
and long-term supply contracts), regulatory diversification
mandates, and political risk-tolerance calibration.  
In other words, we conceptualize the 'planner' as a centralized national authority (e.g., an Energy Security Council) with the mandate to manage systemic risk. While operational flows are executed by private entities, the planner influences the network topology through strategic reserve positioning, infrastructure investment subsidies, and diversification regulations that cap corridor-specific exposure.
%In the Indian case study presented in this paper, these authorities map to an identifiable set of government agencies; Section~\ref{sec:institutional} makes this mapping precise.  The abstraction is transferable: 
For example, the United States (SPR + DoE), China (NDRC + SINOPEC/CNPC coordination), and
the European Union (gas storage obligations under Regulation
2022/1032) all maintain structurally analogous planning functions, differing only in institutional labels.
 
\paragraph{Paper outline}
\Cref{sec:literature} reviews related work. \Cref{sec:problem} describes
the problem setting. \Cref{sec:model} presents the mathematical
formulation, including the decision-dependent probability mechanism and
CVaR integration. \Cref{sec:structural} establishes tractability and
decomposition results.
\Cref{sec:algorithm} develops the Benders decomposition algorithm with
valid inequalities. \Cref{sec:computational} reports computational results.
\Cref{sec:conclusion} concludes the paper. 


%================================================================
\section{Literature Review}
\label{sec:literature}
%================================================================

Our work intersects four streams: stochastic network design under
disruption, risk-averse optimization with CVaR, decision-dependent
(endogenous) uncertainty, and maritime chokepoint and energy-security
resilience.

\paragraph{Stochastic supply chain network design}
The two-stage stochastic programming framework for facility location and
network design under disruption was established by \cite{snyder2005}
and extended to reliable network design by \cite{lim2010}. More recent
contributions include \cite{tolooie2020}, who apply L-shaped
methods for reliable SCND; \cite{fattahi2020}, who
introduce resilience metrics via quadratic conic optimization; and \cite{vaziri2025}, who solve a trilevel
designer--interdictor--designer model for multicommodity network design
under stochastic interdictions via branch-and-Benders-cut. In the maritime
domain, \cite{brouer2014} and \cite{koza2020}
address vessel routing recovery but not strategic chokepoint-level
resilience. Recent procurement-focused work also demonstrates the value of risk-aware stochastic optimization under disruption. \cite{chase2025multi} develop a two-stage multi-period procurement model for COVID-19 lockdown disruptions, combining supplier-risk assessment, supplier diversification, and stock-surplus decisions. Their setting is procurement planning rather than maritime network design, but it reinforces the role of anticipative stochastic optimization in disruption-resilient supply planning.

\paragraph{CVaR in stochastic programming}
\cite{rockafellar2000} showed that CVaR can be optimized via
linear programming through an auxiliary variable reformulation. \cite{noyan2012} developed CVaR-constrained stochastic programs and decomposition
methods for facility location. \cite{carneiro2010risk} applied
mean--CVaR to chemical supply chain disruption mitigation, demonstrating
that CVaR impact increases with decreasing disruption frequency. \cite{liu2020} introduced behavioral CVaR for hazmat network design.
\cite{ghaffarinasab2023} combined mean--CVaR hub location with Benders
decomposition and scenario grouping. \cite{gao2019}
extended the Risk Exposure Index to multiple simultaneous disruptions with
worst-case CVaR, proving that optimal strategic inventory positioning is
characterized by a conic program, though restricted to single-commodity
flow. 


\paragraph{Decision-dependent uncertainty}
\cite{goel2006} distinguish Type~1 (decisions alter
probability distributions) from Type~2 (decisions determine information
timing) endogenous uncertainty. \cite{hellemo2018}
formalize decision-dependent probabilities through affine transformations,
providing the natural formalism for our chokepoint model. \cite{basciftci2021} derive exact MILP reformulations for DRO facility location
with decision-dependent stochastic demand. \cite{bhuiyan2020} model
network design with protection decisions affecting post-disruption capacity
probabilistically. \cite{luo2025two} study supply-chain network design under disruptions and endogenous demand uncertainty using a two-stage stochastic-robust formulation. \cite{zhao2019} develop decision-dependent
adaptive robust optimization for biofuel supply chains. To our knowledge,
no prior paper combines decision-dependent disruption probabilities with
CVaR risk measures in a multi-commodity network design context.

\paragraph{Maritime chokepoint and energy-security resilience}
A separate empirical literature quantifies the systemic exposure of
global trade to maritime chokepoint disruption.
\cite{verschuur2025} use a high-resolution AIS-based simulation
framework to estimate the trade-volume impact of disruptions at major
chokepoints, including Hormuz, Bab~el-Mandeb, and the Suez Canal, but
adopt a descriptive rather than prescriptive stance. The 2024 Houthi
campaign in the Red Sea and the 2026 Strait of Hormuz disruption
prompted policy-oriented analyses by the
\cite{iea2026hormuz}, \cite{cnbc2026pipelines}, and the
\cite{indiastrategic2026hormuz}, documenting Cape rerouting,
war-risk insurance premia, and pipeline-bypass utilization. On the
energy-security side, \cite{eu2022gasstorage} formalizes the European
Union's gas-storage filling mandate and \cite{eia2016spr} document the
expansion of national strategic petroleum reserves. None of these
contributions formulates the joint route-and-inventory design problem
under correlated chokepoint disruption as an optimization model. Our
formulation fills this gap by combining multi-commodity flow with
correlated chokepoint failures, exposure-adjusted probabilities, and
mean--CVaR risk aversion in a prescriptive stochastic MILP.

\paragraph{Position relative to the existing literature}
The literature provides each of the five ingredients of our model in isolation: stochastic network design under disruption \citep{snyder2005, lim2010, tolooie2020, fattahi2020}, multi-commodity interdiction \citep{vaziri2025}, mean--CVaR with decomposition \citep{noyan2012, ghaffarinasab2023}, and decision-dependent probabilities \citep{goel2006, hellemo2018, basciftci2021, bhuiyan2020}. \cite{verschuur2025} provide the closest empirical counterpart in the maritime-chokepoint setting but rely on descriptive simulation rather than prescriptive optimization. Our contribution is to combine all five ingredients—correlated chokepoint failures, decision-dependent probabilities, mean--CVaR risk aversion, commodity-specific admissibility, and corridor-dependence caps—within a single tractable two-stage MILP, and to show that their interaction admits exact reformulation, scenario-decomposable recourse, and finitely convergent Benders decomposition.

\section{Problem Description}
\label{sec:problem}

We consider a directed network $G = (N, A)$ where $N$ is a set of nodes
(supply terminals, transshipment hubs, storage facilities, demand centers)
and $A \subseteq N \times N$ is a set of feasible directed arcs. Multiple
commodities $c \in C$ must be routed from supply nodes $R \subseteq N$ to
demand nodes $D \subseteq N$. Arcs are grouped into \emph{corridors}
$\ell \in \mathcal{L}$ that share a common geographic bottleneck (e.g.,
all arcs transiting the Strait of Hormuz form a single corridor).

The planning problem has two stages:

\begin{enumerate}[leftmargin=2em]
    \item \textbf{First stage (design).} Before uncertainty is resolved,
    the planner determines strategic inventory positions $W_{ic} \ge 0$
    at each node--commodity pair, corresponding to ISPRL cavern
    fills (crude), regasification buffer stocks (LNG), and port
    terminal reserves (LPG, fertilizer), and selects arc activation
    decisions $y_a \in \{0,1\}$ (securing G2G long-term supply
    contracts, sanctioning pipeline bypass infrastructure via
    sovereign guarantees, or chartering dedicated vessel capacity).
    These are ``here-and-now'' decisions.
 

    \item \textbf{Second stage (recourse).} After a disruption scenario
    $s \in S$ is realized---determining which arcs are degraded and by how
    much---the planner routes commodity flows, draws on strategic
    inventory, and incurs shortage penalties for unmet demand.
\end{enumerate}

In the case study, the planner represents the coordinated decision-making apparatus through which a national government manages maritime energy supply-chain risk: strategic reserve operators (e.g., ISPRL in India, the SPR in the United States, NDRC-coordinated entities in China), pipeline-offtake authorities, flag-state regulators, and crisis-response cabinets. The model variables $W_{ic}$, $y_a$, $\bar\Phi_\ell$, $\lambda$, and $\delta_{sa}$ are intentionally generic so that the institutional grounding can be swapped without altering the mathematical structure. \ref{app:institutional} provides the detailed mapping between model variables and the Indian institutional authorities used for calibration.





Three structural features distinguish our model:

\begin{itemize}[leftmargin=2em, itemsep=0.2em]
    \item \textbf{Correlated arc failures.} Arcs within a corridor share
    correlated failure modes. Scenarios $s \in S$ encode joint capacity
    reductions, cost surcharges, transit time increases, and
    commodity-specific admissibility restrictions across all corridors
    simultaneously. A dual-disruption scenario (Hormuz~$+$~Bab el-Mandeb)
    has materially higher probability than the product of marginal
    disruption probabilities.

    \item \textbf{Decision-dependent probabilities.} The probability
    $p_s(\mathbf{y})$ of scenario $s$ depends on first-stage decisions
    $\mathbf{y}$: activating routes through vulnerable corridors shifts
    probability mass toward scenarios in which those corridors are
    disrupted. This formalizes \emph{risk exposure through utilization}.

    \item \textbf{CVaR risk aversion.} The planner minimizes a weighted
    combination of expected recourse cost and the conditional value-at-risk
    at level $\alpha$, protecting against high-cost tail scenarios. In the
    maritime context, this ensures the solution hedges against worst-case
    dual-disruption events, not just expected outcomes.
\end{itemize}


%================================================================
\section{Mathematical Formulation}
\label{sec:model}
%================================================================

We present the formulation in three layers: the base two-stage stochastic
program (\Cref{sec:base-model}), the decision-dependent uncertainty probability mechanism
(\Cref{sec:ddu}), and the CVaR integration (\Cref{sec:cvar}).

%----------------------------------------------------------------
\subsection{Sets, Parameters, and Decision Variables}
\label{sec:notation}
%----------------------------------------------------------------

\begin{table}[H]
\centering
\caption{Sets and indices.}
\label{tab:sets}
\begin{tabular}{@{}ll@{}}
\toprule
Symbol & Description \\
\midrule
$N$ & Set of nodes ($|N|=16$ in the case study) \\
$A \subseteq N \times N$ & Set of feasible directed arcs ($|A|=28$) \\
$C$ & Set of commodities: \{crude, LNG, LPG, fertilizer\} \\
$S$ & Set of disruption scenarios ($|S|=9$) \\
$\mathcal{L}$ & Set of corridors: \{Hormuz, Bab/Suez, Cape, pipeline, direct\} \\
$R \subseteq N$ & Supply nodes (Gulf terminals, non-Gulf suppliers) \\
$D \subseteq N$ & Demand nodes (Indian ports: Jamnagar, Mumbai, etc.) \\
\bottomrule
\end{tabular}
\end{table}

\begin{table}[H]
\centering
\caption{Model parameters.}
\label{tab:params}
\begin{tabularx}{\textwidth}{@{}>{\raggedright\arraybackslash}p{2.5cm}X@{}}
\toprule
Symbol & Description \\
\midrule
$\bar{p}_s$ & Baseline probability of scenario $s$, $\sum_{s} \bar{p}_s = 1$ \\
$\delta_{sa}$ & DDU probability shift: change in $p_s$ per activation of arc $a$ \\
$f_a$ & Fixed activation cost for arc $a$ \\
$c_{ac}^{s}$ & Transport cost on arc $a$, commodity $c$, scenario $s$ \\
$\bar{u}_{ac}^{s}$ & Capacity of arc $a$ for commodity $c$ under scenario $s$ \\
$\theta_{ac}^{s} \in \{0,1\}$ & Admissibility: can commodity $c$ transit arc $a$ in scenario $s$? \\
$d_{ic}$ & Demand for commodity $c$ at node $i \in D$ \\
$r_{ic}$ & Supply of commodity $c$ at node $i \in R$ \\
$L_{ic}$ & Storage capacity for commodity $c$ at node $i$ \\
$h_{c}$ & Holding cost for commodity $c$ per period \\
$\pi_{c}$ & Unmet-demand (shortage) penalty for commodity $c$ \\
$\bar{\Phi}_\ell$ & Corridor-dependence cap for corridor $\ell$ \\
$\lambda \in [0,1]$ & Weight on CVaR term in the mean--CVaR objective \\
$\alpha \in (0,1)$ & CVaR confidence level \\
\bottomrule
\end{tabularx}
\end{table}

\begin{table}[H]
\centering
\caption{Decision variables.}
\label{tab:variables}
\begin{tabularx}{\textwidth}{@{}>{\raggedright\arraybackslash}p{2.5cm}X@{}}
\toprule
Symbol & Description \\
\midrule
\multicolumn{2}{@{}l}{\emph{First-stage (design) variables}} \\
$y_a \in \{0,1\}$ & $= 1$ if arc $a$ is activated \\
$W_{ic} \ge 0$ & Strategic inventory of commodity $c$ pre-positioned at node $i$ \\[0.5em]
\multicolumn{2}{@{}l}{\emph{Second-stage (recourse) variables, indexed by scenario $s$}} \\
$x_{ac}^{s} \ge 0$ & Flow of commodity $c$ on arc $a$, scenario $s$ \\
$I_{ic}^{s} \ge 0$ & Ending inventory of commodity $c$ at node $i$, scenario $s$ \\
$u_{ic}^{s} \ge 0$ & Unmet demand (shortage) for commodity $c$ at node $i$, scenario $s$ \\[0.5em]
\multicolumn{2}{@{}l}{\emph{Risk variables}} \\
$\nu \in \R$ & VaR auxiliary variable \\
$\xi_s \ge 0$ & CVaR excess variable for scenario $s$ \\
\bottomrule
\end{tabularx}
\end{table}


\subsection{Base Two-Stage Stochastic Model}
\label{sec:base-model}

We first present a baseline two-stage stochastic program under two
simplifying assumptions: (i) scenario probabilities are exogenous and
therefore fixed at their baseline values $\bar p_s$, and (ii) the
planner is risk-neutral, so only expected cost is minimized. This base
model provides the benchmark formulation to which we later add
decision-dependent probabilities and mean--CVaR risk aversion.

In the first stage, before the disruption scenario is realized, the
planner chooses arc-activation decisions $y_a$ and strategic inventory
positions $W_{ic}$. In the second stage, after scenario $s \in S$ is
revealed, the planner routes flows through the surviving network,
draws down inventory, and allows unmet demand at a penalty. The second
stage is therefore a recourse problem that evaluates how costly the
first-stage design is under each disruption realization.

\paragraph{Objective}
Under exogenous probabilities and risk neutrality, the objective is
\begin{align}
\min \; Z
&=
\underbrace{\sum_{a \in A} f_a y_a
+ \sum_{i \in N}\sum_{c \in C} h_c W_{ic}}_{\text{first-stage design cost}}
+
\underbrace{\sum_{s \in S} \bar p_s \, Q_s(\mathbf y,\mathbf W)}_{\text{expected second-stage recourse cost}}.
\label{eq:obj-rn}
\end{align}

The first term captures the here-and-now design cost: $f_a y_a$ is the
fixed cost of activating arc $a$, while $h_c W_{ic}$ is the cost of
pre-positioning commodity $c$ at node $i$. The second term is the
expected recourse cost, obtained by weighting the scenario-specific
second-stage cost $Q_s(\mathbf y,\mathbf W)$ by the exogenous scenario
probability $\bar p_s$.

For each scenario $s$, the recourse function is defined as
\begin{align}
Q_s(\mathbf y,\mathbf W)
=
\min \;
& \sum_{a \in A}\sum_{c \in C} c_{ac}^s x_{ac}^s
+ \sum_{i \in N}\sum_{c \in C} h_c I_{ic}^s
+ \sum_{i \in D}\sum_{c \in C} \pi_c u_{ic}^s .
\label{eq:recourse}
\end{align}

The three terms in \eqref{eq:recourse} represent, respectively:
(i) scenario-dependent transportation cost,
(ii) end-of-period inventory holding cost, and
(iii) shortage penalties for unmet demand. Thus, $Q_s(\mathbf y,\mathbf W)$
is the minimum operational cost of responding to disruption scenario $s$
given the first-stage design decisions.

\paragraph{Flow balance}
For every node $i \in N$, commodity $c \in C$, and scenario $s \in S$,
material conservation is enforced by
\begin{align}
W_{ic}
+ r_{ic}
+ \sum_{\arc{j}{i}\in A} x_{ji,c}^{s}
- \sum_{\arc{i}{j}\in A} x_{ij,c}^{s}
- d_{ic}
+ u_{ic}^{s}
=
I_{ic}^{s}.
\label{eq:flow-balance}
\end{align}

Equation \eqref{eq:flow-balance} states that, for each
node--commodity--scenario triple, initial strategic inventory $W_{ic}$,
local supply $r_{ic}$, and inbound flow must cover outbound flow and
demand $d_{ic}$. Any unmet demand is captured by the shortage variable
$u_{ic}^s$, and any remaining amount is recorded as end-of-period
inventory $I_{ic}^s$. Supply and demand are scenario-independent in the
calibrated case study; scenario heterogeneity enters the recourse problem
through arc capacities, costs, and admissibility, not through demand or
supply realizations. This is the core mass-balance equation of the model
and links the network flow decisions to supply, demand, storage, and
shortages.

\paragraph{Capacity constraint}
For each arc $a \in A$, commodity $c \in C$, and scenario $s \in S$,
flow is bounded by scenario-dependent capacity and can occur only if the
arc is activated in the first stage:
\begin{align}
x_{ac}^{s}
\le
\bar u_{ac}^{s}\, y_a,
\quad \forall a \in A,\; c \in C,\; s \in S.
\label{eq:arc-cap}
\end{align}

Constraint \eqref{eq:arc-cap} is the principal linking constraint between
the two stages. If $y_a = 0$, then no recourse flow can use arc $a$.
If $y_a = 1$, then arc $a$ is available up to its scenario-specific
capacity $\bar u_{ac}^{s}$. In this way, first-stage infrastructure or
contracting decisions determine the feasible second-stage response.

% \paragraph{Admissibility.}
% Certain commodities may be prohibited from traversing specific arcs in
% particular scenarios due to physical, regulatory, contractual, or
% geopolitical restrictions. This is modeled by
% \begin{align}
% x_{ac}^{s}
% \le
% \theta_{ac}^{s} M,
% \quad \forall a \in A,\; c \in C,\; s \in S,
% \label{eq:admissibility}
% \end{align}
% where $\theta_{ac}^{s} \in \{0,1\}$ is an admissibility parameter and
% $M$ is a sufficiently large constant. If $\theta_{ac}^{s}=0$, then
% commodity $c$ cannot use arc $a$ in scenario $s$; if $\theta_{ac}^{s}=1$,
% the constraint is nonbinding apart from the arc-capacity limit in
% \eqref{eq:arc-cap}.

%\textcolor{red}{This avoids the use of a loose big-$M$ coefficient.}


\paragraph{Joint Admissibility and Capacity}
To jointly accommodate scenario-specific admissibility and
first-stage arc activation and capacity constraints, we further tighten Constraint \eqref{eq:arc-cap} as follows: 
\begin{align}
x_{ac}^s \le \theta_{ac}^s \bar u_{ac}^s y_a,
\qquad \forall a \in A,\; c \in C,\; s \in S.
\label{eq:admiss-tight}
\end{align}
When $\theta_{ac}^s = 0$, the corresponding flow is forced to zero; when
$\theta_{ac}^s = 1$, the bound reduces to the activated scenario capacity.


\paragraph{Storage capacity}
Strategic and end-of-period inventories are limited by available storage
capacity:
\begin{align}
W_{ic} \le L_{ic}, \qquad
I_{ic}^{s} \le L_{ic},
\quad \forall i \in N,\; c \in C,\; s \in S.
\label{eq:storage}
\end{align}

These bounds ensure that both pre-positioned inventory and residual
inventory remain physically feasible at each node.

\paragraph{Corridor-dependence cap}
To avoid excessive concentration of traffic on a single vulnerable
corridor, we impose a corridor-dependence cap. For each corridor
$\ell \in \mathcal L$ and scenario $s \in S$,
\begin{align}
\sum_{a \in \ell}\sum_{c \in C} x_{ac}^{s}
\le
\bar\Phi_\ell
\sum_{a \in A}\sum_{c \in C} x_{ac}^{s},
\quad \forall \ell \in \mathcal L,\; s \in S.
\label{eq:corridor-cap}
\end{align}

Constraint \eqref{eq:corridor-cap} is a ratio-type diversification
constraint written linearly. It restricts the fraction of total system
flow that can be routed through corridor $\ell$, even if that corridor
is operational and cheap, thereby reducing structural exposure to
concentrated chokepoint risk. The constraint is well-defined whenever
the aggregate scenario flow is strictly positive; in the calibrated
case study, aggregate demand is strictly positive in every scenario and
shortage penalties induce strictly positive total flow at any feasible
recourse solution, so the right-hand side does not collapse. If a
scenario were severe enough that complete network failure forced
$\sum_{a,c} x_{ac}^{s} = 0$, the constraint would be vacuously satisfied
($0 \leq 0$) and could be safely dropped from that scenario's recourse
LP.
%\textcolor{red}{If a scenario is so severe that all demand is unmet ( for all ), the right-hand side becomes zero, so added the below 2 sentences}

Taken together, \eqref{eq:obj-rn}--\eqref{eq:corridor-cap} (less Constraint \eqref{eq:arc-cap} as it is redundant) define the
baseline risk-neutral two-stage stochastic model with fixed scenario
probabilities. In the next subsection, we relax the exogeneity
assumption by allowing scenario probabilities to depend on first-stage
design choices.



\subsection{Decision-Dependent Disruption Probabilities}
\label{sec:ddu}

The baseline model assumes that the disruption probabilities
$\bar p_s$ are exogenous and fixed in advance. In many resilience-planning
settings, however, this assumption is restrictive: the planner's own
first-stage design choices affect how exposed the system is to specific
disruption scenarios. In the present context, a network design that relies
more heavily on vulnerable corridors such as Hormuz or Bab el-Mandeb is
more exposed to scenarios in which those corridors are disrupted. We
therefore allow scenario probabilities to depend on the first-stage arc
activation decisions.

This modeling step should be interpreted as an \emph{exposure-adjusted}
probability mechanism. The idea is not that activating an arc physically
causes a disruption. Rather, a design that depends more strongly on a
fragile corridor shifts the relevant planning probability mass toward
scenarios in which that corridor fails, thereby capturing endogenous
exposure to disruption risk. 
In institutional terms, the $\delta_{sa}$ parameters encode a
forward-looking intelligence assessment produced by MEA, R\&AW,
and the Joint Intelligence Committee, of how each arc activation
shifts the credibility of specific adversarial disruption scenarios,
and are therefore best understood as planning-probability
adjustments rather than physical causal effects.

\begin{definition}[Affine decision-dependent probabilities]
\label{def:ddp}
For each scenario $s \in S$, let the probability of occurrence depend
affinely on the first-stage arc activation vector
$\mathbf y \in \{0,1\}^{|A|}$ according to
\begin{align}
p_s(\mathbf y)
=
\bar p_s + \sum_{a \in A} \delta_{sa}\, y_a,
\qquad \forall s \in S.
\label{eq:ddp}
\end{align}
Here, $\bar p_s$ denotes the baseline probability of scenario $s$, and
$\delta_{sa}$ measures the change in the probability of scenario $s$
induced by activating arc $a$.
\end{definition}



%\textcolor{red}{Your DDP model shifts probabilities based on arc activation, so for this concern I added this}

The dependence of $p_s(\mathbf y)$ on the first-stage activation variables
$y_a$, rather than realized second-stage flows, is intentional. In the present
setting, $y_a$ represents a strategic commitment to a corridor through
contracting, infrastructure activation, or long-term routing dependence. The
resulting probability distortion therefore captures exposure created by the
design itself, rather than by scenario-contingent recourse adjustments. This
choice preserves the two-stage timing structure and keeps the endogenous
uncertainty model computationally tractable.


To ensure that \eqref{eq:ddp} defines a valid probability distribution for
every feasible first-stage decision vector, we impose
\begin{align}
p_s(\mathbf y) \ge 0,
\qquad \forall s \in S,\; \forall \mathbf y \in \{0,1\}^{|A|},
\label{eq:ddp-nonneg}
\end{align}
together with the mass-preservation condition
\begin{align}
\sum_{s \in S} \delta_{sa} = 0,
\qquad \forall a \in A.
\label{eq:ddp-mass}
\end{align}
Condition \eqref{eq:ddp-mass} ensures that activating arc $a$ redistributes
probability mass across scenarios without changing the total mass. Indeed,
summing \eqref{eq:ddp} over all scenarios gives
\[
\sum_{s \in S} p_s(\mathbf y)
=
\sum_{s \in S} \bar p_s
+
\sum_{a \in A} y_a \sum_{s \in S} \delta_{sa}
=
1,
\]
provided that the baseline probabilities satisfy
$\sum_{s \in S} \bar p_s = 1$.

Under \eqref{eq:ddp}, the expected recourse term in the objective becomes
\begin{align}
\sum_{s \in S} p_s(\mathbf y)\, Q_s(\mathbf y,\mathbf W)
&=
\sum_{s \in S}
\left(
\bar p_s + \sum_{a \in A} \delta_{sa} y_a
\right) Q_s(\mathbf y,\mathbf W) \notag \\
&=
\sum_{s \in S} \bar p_s Q_s(\mathbf y,\mathbf W)
+
\sum_{s \in S}\sum_{a \in A} \delta_{sa}\, y_a Q_s(\mathbf y,\mathbf W).
\label{eq:ddp-expand}
\end{align}
Thus, affine decision-dependent probabilities introduce bilinear products
of the form $y_a Q_s$. Analogous bilinear terms arise in the mean--CVaR
extension through products of the form $y_a \xi_s$. Since $y_a$ is binary
and both $Q_s$ and $\xi_s$ are bounded continuous variables, these
products admit an exact mixed-integer linear reformulation via McCormick
linearization, as established in \Cref{prop:milp-reform}.

The affine specification \eqref{eq:ddp} is attractive for two reasons.
First, it is interpretable: the coefficients $\delta_{sa}$ directly encode
how individual arc activations shift exposure toward or away from each
scenario. Second, it preserves tractability, since the resulting
nonlinearities remain binary--continuous bilinear terms that can be
linearized exactly.


\subsection{Mean--CVaR Objective}
\label{sec:cvar}

The baseline model in \Cref{sec:base-model} is risk-neutral: it minimizes
first-stage design cost plus the expected second-stage recourse cost.
While appropriate when average performance is the only concern, such an
objective may underweight rare but severe disruption outcomes. In a
maritime energy-security setting, however, the planner is typically
concerned not only with expected cost, but also with protection against
high-impact tail events such as major chokepoint closures or correlated
multi-corridor disruptions. To capture this tradeoff, we replace the
risk-neutral objective by a mean--CVaR objective.

Specifically, we minimize a convex combination of expected recourse cost
and the conditional value-at-risk (CVaR) of recourse cost. The resulting
objective is
\begin{align}
\min \; Z^{\mathrm{CVaR}}
&=
\sum_{a \in A} f_a y_a
+
\sum_{i \in N}\sum_{c \in C} h_c W_{ic}
+
(1-\lambda)\sum_{s \in S} p_s(\mathbf y)\, Q_s
+
\lambda \left[
\nu + \frac{1}{1-\alpha}\sum_{s \in S} p_s(\mathbf y)\,\xi_s
\right].
\label{eq:obj-cvar}
\end{align}

The first two terms are the deterministic first-stage design costs,
exactly as in the baseline formulation. The third term is the expected
second-stage recourse cost under the scenario distribution
$p_s(\mathbf y)$. The final term is the CVaR component, written in the
standard Rockafellar--Uryasev auxiliary-variable form \cite{rockafellar2000}.

The parameter $\lambda \in [0,1]$ controls the degree of risk aversion.
When $\lambda = 0$, the model reduces to the risk-neutral formulation.
When $\lambda = 1$, the planner optimizes purely with respect to CVaR.
Intermediate values of $\lambda$ generate a tradeoff between average
performance and tail-risk protection.

The parameter $\alpha \in (0,1)$ is the CVaR confidence level. Larger
values of $\alpha$ place greater emphasis on the worst $(1-\alpha)$
fraction of outcomes. For example, $\alpha = 0.9$ means that the model
protects against the average cost incurred in the worst 10\% of the
scenario distribution.

To express CVaR within a tractable optimization model, we introduce the
auxiliary variable $\nu \in \mathbb R$, which plays the role of the
value-at-risk threshold, together with nonnegative excess variables
$\xi_s$ for each scenario $s \in S$. These are linked by
\begin{align}
\xi_s &\ge Q_s - \nu, && \forall s \in S,
\label{eq:cvar-excess} \\
\xi_s &\ge 0, && \forall s \in S.
\label{eq:cvar-nonneg}
\end{align}

At optimality, \eqref{eq:cvar-excess}--\eqref{eq:cvar-nonneg} imply
\[
\xi_s = \max\{Q_s-\nu,0\},
\]
so $\xi_s$ measures the amount by which scenario-$s$ recourse cost
exceeds the threshold $\nu$. Consequently, the quantity
\[
\nu + \frac{1}{1-\alpha}\sum_{s \in S} p_s(\mathbf y)\,\xi_s
\]
is exactly the CVaR of the recourse cost distribution at confidence level
$\alpha$. This representation is especially attractive because it
preserves linearity once the decision-dependent probability terms are
handled through the McCormick reformulation described in
\Cref{prop:milp-reform}.

The mean--CVaR formulation therefore yields a planner-facing objective
that balances operational efficiency against resilience to tail events.
In the empirical study, we set $\lambda = 0.5$ and $\alpha = 0.95$,
corresponding to a moderately risk-averse planner who places equal
weight on expected performance and 95\%-tail protection.


\section{Tractability and Decomposition Structure}
\label{sec:structural}


We establish three results that underpin the solution algorithm (which will be presented in the next section). The
first reformulates the bilinear interaction between decision-dependent
probabilities and the recourse cost as an exact MILP via
binary--continuous McCormick linearization. The second establishes
that, with the McCormick reformulation in place, the recourse problem
remains separable across scenarios and the CVaR component does not
break this separability. The third uses both to deduce finite
convergence of the Benders decomposition that exploits the separable
structure. The three results are tractability statements rather than
deep structural theorems; their role is to certify that the
five-feature integration introduced in \Cref{sec:model} can be solved
exactly without resorting to heuristic decomposition.

\begin{proposition}[Exact MILP reformulation under affine decision-dependent probabilities]
\label{prop:milp-reform}
Under the affine decision-dependent probability model
\begin{align}
p_s(\mathbf{y}) = \bar p_s + \sum_{a \in A} \delta_{sa} y_a,
\qquad \forall s \in S,
\label{eq:ddp-recall}
\end{align}
the mean--CVaR objective \eqref{eq:obj-cvar} admits an exact mixed-integer
linear reformulation.

In particular, introduce auxiliary variables
\[
w_{as} = y_a Q_s, \qquad v_{as} = y_a \xi_s,
\qquad \forall a \in A,\; s \in S,
\]
and suppose valid upper bounds
\[
0 \le Q_s \le M_{as}^Q, \qquad 0 \le \xi_s \le M_{as}^{\xi}
\]
are available. Then the bilinear products can be represented exactly by
the McCormick inequalities
\begin{align}
w_{as} &\ge Q_s - M_{as}^Q(1-y_a), &
v_{as} &\ge \xi_s - M_{as}^{\xi}(1-y_a),
\label{eq:mc-lb} \\
w_{as} &\le Q_s, &
v_{as} &\le \xi_s,
\label{eq:mc-ub1} \\
w_{as} &\le M_{as}^Q y_a, &
v_{as} &\le M_{as}^{\xi} y_a,
\label{eq:mc-ub2} \\
w_{as} &\ge 0, &
v_{as} &\ge 0.
\label{eq:mc-nn}
\end{align}
Consequently, the full objective can be written as a linear function of
\[
(\mathbf y,\mathbf W,\mathbf Q,\boldsymbol{\xi},
\mathbf w,\mathbf v,\nu).
\]
\end{proposition}

\begin{proof}
Substituting \eqref{eq:ddp-recall} into the mean--CVaR objective
\eqref{eq:obj-cvar} yields terms of the form
\[
y_a Q_s
\quad\text{and}\quad
y_a \xi_s,
\]
arising from the products $p_s(\mathbf y)Q_s$ and
$p_s(\mathbf y)\xi_s$. Since $y_a \in \{0,1\}$ is binary and both
$Q_s$ and $\xi_s$ are bounded continuous variables, each such product
admits an exact McCormick linearization. The constraints
\eqref{eq:mc-lb}--\eqref{eq:mc-nn} enforce
\[
w_{as}=y_aQ_s, \qquad v_{as}=y_a\xi_s
\]
exactly. Replacing every occurrence of $y_aQ_s$ and $y_a\xi_s$ in the
objective by $w_{as}$ and $v_{as}$ therefore yields an equivalent MILP
reformulation. 
\end{proof}

\begin{proposition}[Scenario-wise decomposability under mean--CVaR]
\label{prop:cvar-decomp}
For fixed first-stage decisions $(\mathbf{y}, \mathbf{W}, \nu)$, the
second-stage problem remains separable across scenarios. In particular,
for each scenario $s \in S$, the associated recourse variables and CVaR
excess variable $\xi_s$ can be optimized independently of all other
scenarios.
\end{proposition}

\begin{proof}
Fix $(\mathbf{y}, \mathbf{W}, \nu)$. Then the decision-dependent
probabilities $p_s(\mathbf{y})$ are constants, and the first-stage
variables no longer appear as optimization variables in the recourse
stage.

The uncertain part of the objective is
\[
(1-\lambda)\sum_{s\in S} p_s(\mathbf{y}) Q_s
\;+\;
\frac{\lambda}{1-\alpha}\sum_{s\in S} p_s(\mathbf{y}) \xi_s,
\]
since the term $\lambda \nu$ is constant once $\nu$ is fixed. For each
scenario $s$, the CVaR constraints are
\[
\xi_s \ge Q_s - \nu, \qquad \xi_s \ge 0.
\]
Hence, at optimality,
\[
\xi_s = \max\{Q_s-\nu,0\},
\]
so the contribution of scenario $s$ to the objective is
\[
(1-\lambda)p_s(\mathbf{y})Q_s
\;+\;
\frac{\lambda}{1-\alpha}p_s(\mathbf{y})\max\{Q_s-\nu,0\}.
\]

Moreover, the recourse constraints for scenario $s$ (flow balance,
capacity, admissibility, storage, and corridor-cap constraints) involve
only scenario-$s$ variables. No recourse variable from scenario $s$ appears
in the constraints of any other scenario $s' \neq s$.

Therefore, once $(\mathbf{y}, \mathbf{W}, \nu)$ is fixed, the second-stage
optimization decomposes into independent scenario subproblems, one for
each $s \in S$. The only cross-scenario coupling is through the shared
first-stage decisions, which are handled in the master problem.
\end{proof}

\begin{proposition}[Finite convergence of Benders decomposition]
\label{prop:benders_finite}
The Benders decomposition algorithm described in \Cref{sec:algorithm},
using optimality cuts generated from feasible scenario subproblems and
feasibility cuts generated from infeasible scenario subproblems,
terminates after finitely many iterations and returns an optimal solution
of the extensive-form model.
\end{proposition}

\begin{proof}[Proof sketch]
By \Cref{prop:cvar-decomp}, for fixed first-stage decisions
$(\mathbf{y},\mathbf{W},\nu)$, the second-stage problem decomposes into
independent scenario subproblems. Each such subproblem is a linear
program. Hence, its dual feasible region is a polyhedron with finitely
many extreme points and extreme rays.

Feasible scenario subproblems generate valid Benders optimality cuts,
while infeasible scenario subproblems generate valid feasibility cuts.
Each cut is satisfied by every feasible solution of the extensive form,
so the master problem remains a relaxation of the extensive-form model
throughout the algorithm.

At each iteration, the current master solution is evaluated by solving
the scenario subproblems. If all subproblems are feasible and the
current recourse approximation is exact, then the current master
solution is optimal for the extensive form and the algorithm terminates.
Otherwise, at least one violated optimality or feasibility cut is
identified and added to the master problem. Since the number of distinct
cuts induced by the finitely many dual extreme points and rays is finite,
this process cannot continue indefinitely.

Therefore, after finitely many iterations, no further violated cut
exists, and the master solution coincides with an optimal solution of
the extensive-form model.
\end{proof}


%================================================================
\section{Solution Algorithm}
\label{sec:algorithm}
%================================================================

The extensive-form model can be solved directly on the case-study
instances, but its structure is naturally amenable to decomposition.
We therefore solve the model by Benders decomposition, exploiting the
scenario-wise separability established in
\Cref{prop:cvar-decomp}. The key idea is to keep the first-stage design
variables in a master problem and evaluate the true second-stage
recourse cost through independent scenario subproblems.

Because the scenario probabilities $p_s(\mathbf y)$ are
decision-dependent, the master objective contains bilinear products such
as $p_s(\mathbf y)\phi_s$ and $p_s(\mathbf y)\eta_s$. As in
\Cref{prop:milp-reform}, these terms are handled exactly through
McCormick linearization, so the resulting master problem becomes a mixed-integer linear
program.

%----------------------------------------------------------------
\subsection{Master Problem}
\label{sec:master}
%----------------------------------------------------------------

The master problem determines the first-stage design decisions
$(\mathbf y,\mathbf W,\nu)$ together with scenario-wise recourse
underestimators $\phi_s \ge 0$ and CVaR excess variables
$\eta_s \ge 0$, for all $s \in S$:
\begin{align}
\min \;
& \sum_{a \in A} f_a y_a
+ \sum_{i \in N}\sum_{c \in C} h_c W_{ic}
+ (1-\lambda)\sum_{s \in S} p_s(\mathbf y)\,\phi_s
+ \lambda\left[
\nu + \frac{1}{1-\alpha}\sum_{s \in S} p_s(\mathbf y)\,\eta_s
\right]
\label{eq:master-obj}
\end{align}
subject to
\begin{align}
\eta_s &\ge \phi_s - \nu,
&& \forall s \in S,
\label{eq:master-cvar-link} \\
\eta_s &\ge 0,
&& \forall s \in S,
\label{eq:master-cvar-nonneg}
\end{align}
together with the first-stage domain constraints, storage-capacity
constraints, probability-validity conditions for $p_s(\mathbf y)$, and
all Benders cuts accumulated up to the current iteration.

The variable $\phi_s$ is a master-level lower approximation of the true
scenario recourse cost $Q_s(\mathbf y,\mathbf W)$. The variable
$\eta_s$ plays the corresponding role for the CVaR excess term. At
convergence, the master satisfies
\[
\phi_s = Q_s(\mathbf y,\mathbf W),
\qquad
\eta_s = \max\{Q_s(\mathbf y,\mathbf W)-\nu,0\},
\qquad \forall s \in S.
\]

%----------------------------------------------------------------
\subsection{Scenario Subproblems}
\label{sec:subproblem}
%----------------------------------------------------------------

Given a candidate first-stage solution
$(\hat{\mathbf y},\hat{\mathbf W},\hat\nu)$ from the master, the
scenario-$s$ subproblem is the recourse LP
\begin{align}
Q_s(\hat{\mathbf y},\hat{\mathbf W})
=
\min \;
& \sum_{a \in A}\sum_{c \in C} c_{ac}^s x_{ac}^s
+ \sum_{i \in N}\sum_{c \in C} h_c I_{ic}^s
+ \sum_{i \in D}\sum_{c \in C} \pi_c u_{ic}^s
\label{eq:sub-obj}
\end{align}
subject to the scenario-$s$ recourse constraints:
flow balance, arc-capacity linking, admissibility, storage-capacity,
and corridor-dependence constraints, with
$(\mathbf y,\mathbf W)$ fixed at
$(\hat{\mathbf y},\hat{\mathbf W})$.

Since the first-stage decisions are fixed, each scenario subproblem is a
linear program. Its dual optimal solution yields a Benders optimality
cut that lower-bounds the recourse function in the master problem. In
generic form, such a cut can be written as
\begin{align}
\phi_s
\;\ge\;
\Theta_s
+ \sum_{a \in A} \Gamma_{as} y_a
+ \sum_{i \in N}\sum_{c \in C} \Lambda_{ics} W_{ic},
\label{eq:benders-opt-cut-generic}
\end{align}
where the coefficients
$(\Theta_s,\Gamma_{as},\Lambda_{ics})$ are obtained from the dual
multipliers of the scenario-$s$ LP. If the scenario subproblem is
infeasible, then a dual extreme ray generates a Benders feasibility cut,
which is added to the master problem to exclude the current first-stage
solution.

%----------------------------------------------------------------
\subsection{Corridor-Based Strengthening Cuts}
\label{sec:valid-ineq}
%----------------------------------------------------------------

Standard Benders cuts are generated independently by scenario. We
further strengthen the master relaxation by exploiting the
corridor-based failure structure of the scenario set.

\begin{remark}[Group-failure valid inequality]
\label{prop:group-cut}
Let $\ell \in \mathcal L$ be a corridor, and let
$S_\ell \subseteq S$ denote the set of scenarios in which corridor
$\ell$ is disrupted as a group. Then, for each $s \in S_\ell$, the
master problem can be strengthened by the valid inequality
\begin{align}
\phi_s \ge \underline Q_s^\ell
+ \sum_{a \in \ell} \mu_a^\ell (1-y_a),
\qquad \forall s \in S_\ell,
\label{eq:group-cut}
\end{align}
where $\underline Q_s^\ell$ is a valid lower bound on the recourse cost
in scenario $s$ when all arcs in corridor $\ell$ are activated, and
$\mu_a^\ell \ge 0$ is the marginal increase in recourse cost caused by
deactivating arc $a$ under corridor-$\ell$ disruption.
\end{remark}

These cuts exploit structural information shared across scenarios that
involve common corridor failures. In practice, they tighten the master
problem and can reduce the number of Benders iterations required for
convergence.


\subsection{Algorithm Summary}
\label{sec:alg-summary}


\begin{algorithm}[H]
\caption{Benders decomposition with corridor-based strengthening cuts}
\label{alg:benders}
\begin{algorithmic}[1]
\State \textbf{Initialize:}
LB $\gets -\infty$, UB $\gets +\infty$, tolerance $\epsilon > 0$.
\Repeat
    \State Solve the current master problem
    \eqref{eq:master-obj}--\eqref{eq:master-cvar-nonneg}
    and obtain
    $(\hat{\mathbf y},\hat{\mathbf W},\hat\nu,
    \hat{\boldsymbol\phi},\hat{\boldsymbol\eta})$;
    update the lower bound LB.
    \For{each scenario $s \in S$}
        \State Solve the scenario-$s$ subproblem
        \eqref{eq:sub-obj} with fixed
        $(\hat{\mathbf y},\hat{\mathbf W})$.
        \If{the subproblem is feasible}
            \State Obtain the optimal value $Q_s^*$ and the associated dual
            solution.
            \If{$Q_s^* > \hat\phi_s + \epsilon$}
                \State Add the corresponding Benders optimality cut
                \eqref{eq:benders-opt-cut-generic} to the master.
            \EndIf
        \Else
            \State Generate a Benders feasibility cut from a dual extreme
            ray and add it to the master.
        \EndIf
    \EndFor
    \State Generate violated corridor-based strengthening cuts
    \eqref{eq:group-cut} and add them to the master.
    \State Evaluate the current first-stage solution using the true
    scenario recourse values and update the upper bound UB.
\Until{$(\mathrm{UB}-\mathrm{LB})/\max\{1,|\mathrm{UB}|\} < \epsilon$}
\end{algorithmic}
\end{algorithm}



\section{Computational Study}
\label{sec:computational}

The computational study below uses a 16-node, 28-arc network calibrated
to Indian maritime energy imports, with four commodities (crude oil,
LNG, LPG, fertilizer) and a nine-scenario disruption library. 
Note that this may sound like a relatively small instance, but what we are interested to study is a "national-level" abstraction where the 16 nodes represent major aggregate hubs rather than every single port.
The national
import anchors (5.5\,mb/d crude, 17.5\,Mt/yr LNG, 23.3\,Mt/yr LPG,
9.2\,Mt/yr fertilizer feedstock) are sourced from EIA, UNCTAD, PPAC,
and Drewry; ISPRL strategic-storage capacities are sourced from PIB
\citep{pib2021spr}; pipeline-bypass capacities (ADCOP 1.5\,mb/d,
Petroline 7\,mb/d) are sourced from \cite{iea2026hormuz} and
\cite{cnbc2026pipelines}; and disruption multipliers are anchored to
UKMTO advisories, EIA Red-Sea-decline analyses, and UNCTAD
post-Houthi-crisis reporting. The full calibration protocol, with each
parameter's source, anchor, and construction rule, is documented in
\ref{app:calibration}; all numerical inputs are stored in CSV
files alongside the implementation, allowing alternative calibrations
without code changes. Risk parameters are set at $\lambda=0.5$ and
$\alpha=0.95$ in the base case and swept in
\Cref{sec:cvar-frontier}. The model is implemented in Pyomo and solved
with GLPK; all experiments run on a single thread of a workstation-
class CPU, with extensive-form solve times under one second on the
calibrated instance.

\subsection{Base-Case Experiment: Full Resilience Model}
\label{sec:base-case}

We begin with a base-case experiment that represents the planner's
current energy-security problem under the full proposed model. The
setting is an import-dependent maritime energy network in which crude
oil, LNG, LPG, and fertilizer feedstock must be delivered to Indian
demand nodes through a set of maritime, pipeline, and direct supply
arcs. Several of these arcs pass through geopolitically sensitive
corridors, including Hormuz and Bab/Suez, whose disruptions may occur
jointly rather than independently.

Before the disruption scenario is realized, the planner chooses which
arcs to activate and how much strategic inventory to pre-position at
eligible nodes. After a scenario is realized, the planner routes
available commodity flows through the surviving network, carries
remaining inventory, and incurs shortage penalties for unmet demand.
Thus, the base-case solution represents a complete here-and-now
resilience policy and its scenario-contingent operational response.

The base-case model includes all three structural features introduced
in the paper: correlated disruption scenarios, affine
decision-dependent probabilities, and mean--CVaR risk aversion. Unless
otherwise stated, we use $\lambda=0.5$ and $\alpha=0.95$, corresponding
to a planner that places equal weight on expected recourse cost and
95\%-tail protection. This experiment establishes the reference design
against which the value-of-modeling and algorithmic experiments are
compared.

\paragraph{First-stage network design}
Figure~\ref{fig:base-network-design} visualizes the optimized
first-stage network design. The highlighted arcs correspond to
activated routes, contracts, or bypass options selected before the
disruption scenario is known. The selected design activates 21 of the
28 candidate arcs, indicating that the planner does not rely only on a
single low-cost maritime corridor. Instead, the solution builds a
diversified route portfolio containing Gulf routes, direct routes,
pipeline-bypass links, and Cape alternatives. This is consistent with
the role of the first-stage variable $y_a$: it represents a strategic
commitment to route availability before the uncertain disruption state
is realized.

\begin{figure}[t]
    \centering
    \includegraphics[width=\textwidth]{plots/fig01_base_network_design.pdf}
    \caption{Base-case first-stage network design. Highlighted arcs indicate
    routes activated before the disruption scenario is realized. Node markers
    distinguish supply terminals, transshipment or bypass nodes, and Indian
    demand ports.}
    \label{fig:base-network-design}
\end{figure}

\paragraph{Cost decomposition}
Table~\ref{tab:tab01_base_case} reports the cost decomposition of the
base-case solution. The optimal mean--CVaR objective is 59{,}494.77,
with a first-stage design cost of 1{,}552.00 and an expected recourse
cost of 45{,}673.67. The optimized CVaR$_{0.95}$ value is 70{,}211.88,
which is equal to the VaR threshold in this discrete nine-scenario
instance. Hence, for the calibrated base case, the tail-risk term is
governed by the worst realized scenario in the scenario library. The
extensive form solves in 0.03 seconds, so this instance serves as a
transparent reference case for interpreting the structure of the
solution before moving to the larger value-of-modeling and
decomposition experiments.

\begin{table}
\caption{Base-case solution headline metrics.}
\label{tab:tab01_base_case}
\begin{tabular}{rrrrrrr}
\toprule
objective & fixed\_cost & expected\_recourse & cvar\_0.95 & var\_threshold & n\_arcs\_activated & runtime\_sec \\
\midrule
59494.77 & 1552.00 & 45673.67 & 70211.88 & 70211.88 & 21 & 0.03 \\
\bottomrule
\end{tabular}
\end{table}


\paragraph{Scenario-contingent recourse and rerouting}
The first-stage activation decisions are fixed across all scenarios,
but the second-stage flows are scenario-dependent. Figure~\ref{fig:basecase-rerouting}
illustrates this distinction by comparing the normal operating pattern
with the flow pattern under a disruption scenario. In the normal
scenario, the model routes substantial flow through the conventional
Gulf-to-India corridors and direct supply routes. Under disruption,
the same activated network is used differently: flows shift away from
the disrupted chokepoint and toward direct and Cape alternatives where
capacity remains available. Thus, the ``after-disruption'' network
should be interpreted as a recourse-flow pattern $x^s_{ac}$, not as a
new activation decision.

\begin{figure}[t]
    \centering
    \includegraphics[width=\textwidth]{plots/fig01_before_after_disruption_flows.pdf}
    \caption{Scenario-contingent recourse flows under the base-case design.
    The left panel shows flows in the normal scenario, while the right panel
    shows flows under the selected disruption scenario. Line width is
    proportional to total arc flow aggregated over commodities.}
    \label{fig:basecase-rerouting}
\end{figure}

The corridor-flow decomposition confirms the same behavior numerically.
In the normal scenario, the solution uses Hormuz and direct routes as
the dominant flow channels. When Hormuz capacity is reduced or closed,
the flow portfolio changes sharply: direct routes and Cape alternatives
absorb a much larger share of system flow, while Hormuz flow either
falls substantially or becomes zero.

\paragraph{Scenario-level losses}
Figure~\ref{fig:scenario-costs} reports scenario-level recourse costs
under the base-case design. The results show that disruption severity
is highly asymmetric across scenarios. Several scenarios can be handled
mainly through rerouting and inventory use, while Hormuz-related
disruptions generate substantial residual shortage penalties. The
highest-cost case is the dual-disruption scenario, with recourse cost
70{,}211.88. Hormuz closure and closure-with-bypass scenarios also
produce large shortage costs. These outcomes indicate that the dominant
residual risk is not generic transport-cost inflation alone, but the
combination of reduced Hormuz capacity, commodity-specific admissibility
restrictions, and limited substitutability of certain energy
commodities.

\begin{figure}[t]
    \centering
    \includegraphics[width=\textwidth]{plots/fig01_scenario_costs.pdf}
    \caption{Scenario-level recourse costs under the base-case first-stage
    design.}
    \label{fig:scenario-costs}
\end{figure}

\paragraph{Strategic inventory and commodity-level vulnerability}
Table~\ref{tab:tab01_commodity_outcomes} summarizes the commodity-level
inventory and shortage outcomes. The solution fills available storage
for crude oil, LNG, and LPG, and uses most of the available fertilizer
storage. Crude oil and fertilizer experience no shortage across the
calibrated scenario library, indicating that the selected combination
of strategic reserves and routing alternatives is sufficient to hedge
these commodities in the base case. In contrast, LNG and LPG remain
vulnerable. LPG has the largest expected shortage quantity and the
largest expected shortage cost, identifying it as the most exposed
commodity under the optimized base-case policy. This finding is
consistent with the operational interpretation that LPG has fewer
effective substitution and storage options than crude oil.

Quantitatively, expected LPG shortage (200.84 units, 10{,}042 in expected
shortage cost) is roughly twice the LNG figure (109.31 units, 4{,}372 in
cost) and is concentrated in scenarios that combine Hormuz capacity
reduction with Bab/Suez disruption; crude oil and fertilizer record zero
expected and zero worst-case shortage across the calibrated scenario
library.

\begin{table}[ht]
\centering
\caption{Commodity-level shortage and inventory outcomes.}
\label{tab:tab01_commodity_outcomes}
\small % Slightly smaller font helps with dense data
\begin{tabularx}{\textwidth}{lXXXXXX}
\toprule
\textbf{Commodity} & \textbf{Exp. Shortage Qty} & \textbf{Exp. Shortage Cost} & \textbf{Worst Case Qty} & \textbf{Total Inv.} & \textbf{Storage Cap.} & \textbf{Inv. Util.} \\
\midrule
Crude      & 0.00   & 0.00     & 0.00   & 2700.00 & 2700.00 & 1.00 \\
LNG        & 109.31 & 4372.48  & 250.00 & 450.00  & 450.00  & 1.00 \\
LPG        & 200.84 & 10042.00 & 350.00 & 400.00  & 400.00  & 1.00 \\
Fertilizer & 0.00   & 0.00     & 0.00   & 700.00  & 800.00  & 0.88 \\
\bottomrule
\end{tabularx}
\end{table}

The optimized policy places crude reserves at major refining or coastal demand nodes, LNG buffers at regasification-linked nodes, and LPG and fertilizer inventories at selected port nodes. These inventory decisions are part of the first-stage design and are therefore chosen before any disruption is observed.

\paragraph{Decision-dependent probabilities in the base case}
For this particular optimized design, the DDU-adjusted probabilities
are numerically close to the baseline probabilities. This should not be
interpreted as removing the relevance of the decision-dependent
probability model. Rather, it means that the selected arc portfolio
induces offsetting probability shifts across scenarios in the calibrated
base case. The role of endogenous probabilities is therefore isolated
more directly in the value-of-DDU experiment, where a design optimized
without decision-dependent probabilities is evaluated under the true
DDU-adjusted distribution.

\paragraph{Takeaway}
Overall, the base-case solution shows that the proposed model produces
an interpretable resilience policy. Before disruption, the planner
activates a diversified set of routes and fills most available
strategic storage. After disruption, the recourse solution reroutes
flows away from impaired chokepoints toward direct and Cape alternatives
where possible. At the commodity level, crude oil and fertilizer are
fully protected in the calibrated scenario library, whereas LNG and
especially LPG retain residual shortage exposure. These results motivate
the subsequent experiments, which isolate the value of stochastic
modeling, correlated disruptions, endogenous probabilities, and
risk-averse planning.








\subsection{Value of Stochastic Modeling}
\label{sec:vss}

The base-case experiment evaluates the full resilience model under the
calibrated disruption distribution. We next ask whether explicitly
modeling scenario uncertainty changes the quality of the first-stage
design. To isolate this effect, we temporarily disable
decision-dependent probabilities and CVaR risk aversion and compare
three risk-neutral planning models.

The first model is the stochastic program (SP), which optimizes a
single first-stage design under the full scenario distribution. The
second is the expected-value (EV) approximation, which replaces the
scenario library by a single probability-weighted average scenario,
solves the resulting deterministic model, and then evaluates the
resulting first-stage design under the original stochastic scenario
distribution. The third is the wait-and-see (WS) benchmark, which solves
a separate deterministic problem for each scenario as if the planner had
perfect advance information. The WS solution is not implementable, but
it provides a perfect-information lower bound.

We define the value of the stochastic solution as
\[
\mathrm{VSS}
=
Z_{\mathrm{EV}}
-
Z_{\mathrm{SP}},
\]
where \(Z_{\mathrm{EV}}\) denotes the cost of the expected-value design
when evaluated under the full scenario distribution. We also report the
expected value of perfect information,
\[
\mathrm{EVPI}
=
Z_{\mathrm{SP}}
-
Z_{\mathrm{WS}}.
\]
A positive VSS indicates that optimizing against an average disruption
state leads to a first-stage design that is suboptimal under the true
distribution of disruption outcomes.



\begin{table}[ht]
\centering
\caption{Value of the stochastic solution and EVPI.}
\label{tab:tab02_voss}
\renewcommand{\theadfont}{\bfseries\small} % Makes headers bold and smaller
\begin{tabular}{ccccccc}
\toprule
\thead{SP\\Objective} & \thead{EV Design\\Full Dist.} & \thead{WS Perfect\\Info} & \thead{VSS} & \thead{VSS\\(\%)} & \thead{EVPI} & \thead{EVPI\\(\%)} \\
\midrule
47417.02 & 54446.89 & 47406.64 & 7029.87 & 14.83 & 10.38 & 0.02 \\
\bottomrule
\end{tabular}
\end{table}

Table~\ref{tab:tab02_voss} shows that the stochastic-programming
design has objective value \(47{,}417.02\), while the expected-value
design, when evaluated under the full scenario distribution, has
objective value \(54{,}446.89\). The resulting value of the stochastic
solution is \(7{,}029.87\), corresponding to a 14.83\% increase relative
to the stochastic-program objective. This gap is operationally
significant: it shows that planning against a single average disruption
state produces a materially worse first-stage design once the true
scenario distribution is restored.

In contrast, the wait-and-see benchmark has objective value
\(47{,}406.64\), only \(10.38\) below the stochastic-program objective.
The resulting EVPI is 0.02\% of the stochastic-program objective. This
small EVPI indicates that, once the stochastic program is used, the
remaining value of perfect advance information is limited. The expected ordering $Z_{\mathrm{WS}}\le Z_{\mathrm{SP}}\le Z_{\mathrm{EV}}$ holds, with the near equality between WS and SP indicating that the stochastic program is close to the perfect-information lower bound, while the much larger EV objective shows that the deterministic average-scenario approximation fails to preserve the relevant disruption structure. The main value therefore comes from modeling the distribution of possible disruption states, not from perfect ex-post information.

\begin{figure}[t]
    \centering
    \includegraphics[width=0.82\textwidth]{plots/fig02_design_difference.pdf}
    \caption{Comparison of first-stage route activations by corridor under
    the stochastic-program and expected-value designs.}
    \label{fig:vss-design-difference}
\end{figure}

Figure~\ref{fig:vss-design-difference} explains the source of the VSS
gap. The stochastic-program and expected-value designs do not activate
the same route portfolio. The EV design is optimized for an averaged
network state and therefore under-represents the value of route
redundancy across distinct disruption realizations. The SP design, by
contrast, internalizes the full distribution of scenario-contingent
capacity and cost changes, leading to a first-stage policy that performs
better when evaluated across the actual scenario library.

Overall, this experiment confirms that stochastic modeling is not a
cosmetic extension of a deterministic network design model. The
expected-value approximation produces a substantially worse
implementable policy, while the stochastic program nearly closes the
gap to the perfect-information benchmark. This justifies the use of a
scenario-based stochastic formulation before introducing the additional
features of correlated disruptions, endogenous probabilities, and
risk-averse CVaR planning.


\subsection{Value of Modeling Correlated Chokepoint Disruptions}
\label{sec:vmc}
 
Experiment~\ref{sec:vss} showed that optimizing against the full
scenario distribution produces a substantially better implementable
design than optimizing against a single expected-value scenario. We now
ask whether the dependence structure inside the scenario distribution
also changes the first-stage policy. This is relevant because maritime
energy disruptions are not necessarily independent: a regional
escalation can simultaneously affect Hormuz, Bab/Suez, insurance
conditions, rerouting costs, and tanker availability.
 
To isolate this effect, we compare the calibrated correlated scenario
distribution with an independent-marginal counterfactual. The
counterfactual preserves the marginal disruption probabilities of the
main corridors but removes the joint dependence structure among them.
Both models are solved using a risk-neutral fixed-probability
configuration, with decision-dependent probabilities and CVaR disabled.
The design optimized under the independent-marginal distribution is
then evaluated under the original correlated distribution.
 
Let \(Z_{\mathrm{Corr}}\) denote the optimal objective under the
correlated scenario library, and let \(Z_{\mathrm{Ind}\rightarrow
\mathrm{Corr}}\) denote the objective obtained by evaluating the
independence-based design under the correlated distribution. We define
the value of modeling correlation as
\[
\mathrm{VMC}
=
Z_{\mathrm{Ind}\rightarrow\mathrm{Corr}}
-
Z_{\mathrm{Corr}}.
\]
A positive value indicates that ignoring joint chokepoint dependence
leads to a first-stage design that performs worse under the true
correlated disruption environment.
 
\begin{table}
\caption{Value of modeling correlation.}
\label{tab:tab03_vmc}
\begin{tabular}{rrrrrr}
\toprule
Corr\_objective & Ind\_objective & Ind\_design\_under\_corr & VMC & VMC\_pct & ordering\_holds \\
\midrule
47417.02 & 48094.16 & 47417.02 & 0.00 & 0.00 & True \\
\bottomrule
\end{tabular}
\end{table}

 
Table~\ref{tab:tab03_vmc} reports the correlation experiment. The
correlated model has objective value \(47{,}417.02\). The
independent-marginal model has a different in-sample objective,
\(48{,}094.16\), under its own probability distribution. However, when
the independent design is evaluated under the correlated distribution,
its objective is \(47{,}417.02\), identical to the correlated optimum.
Therefore, the measured VMC is zero for the calibrated base instance.
 
\paragraph{Structural joint-failure resilience}
 
We investigate whether the zero VMC is a calibration artifact or a
structural property. Two parametric experiments are conducted. First, we
interpolate between the independent and correlated distributions using a
parameter \(\rho\in[0,1]\), where \(p(\rho)=(1-\rho)\,p_{\mathrm{ind}}
+\rho\,p_{\mathrm{corr}}\), and optimize under \(p(\rho)\) for each
\(\rho\in\{0,0.1,\ldots,1\}\). The resulting design is then evaluated
under the true correlated distribution. Across all eleven values of
\(\rho\), the VMC remains exactly zero: every mis-specified design
evaluates to the same objective as the correlated optimum. The KL
divergence between the planning and true distributions grows to 0.026
at \(\rho=0\), confirming that the probability distributions are
meaningfully different even though the optimal design is invariant.
 
Second, we amplify the probability of joint-disruption scenarios by
factors \(\gamma\in\{1,1.5,2,3,5,8,10,15\}\) and renormalize. At
\(\gamma=15\), the joint-failure probability (Hormuz and Bab/Suez
simultaneously disrupted) reaches 0.87 under the correlated
distribution versus 0.84 under independence—a nontrivial
divergence. Nevertheless, the VMC remains exactly zero across all
amplification levels, and the optimal design is identical (13 activated
arcs) under both the correlated and independent distributions at every
tested amplification.
 
\begin{figure}[t]
    \centering
    \includegraphics[width=0.95\textwidth]{plots/fig13_stress_test.pdf}
    \caption{VMC under amplified joint-failure probabilities. Left: VMC
    remains zero across all amplification levels. Right: the gap between
    correlated and independent joint-failure probabilities grows with
    amplification, yet the optimal design does not change.}
    \label{fig:vmc-stress}
\end{figure}
 
\begin{table}[htbp] % Standard table environment instead of sidewaystable
\centering
\caption{VMC under amplified joint-failure correlation.}
\label{tab:tab13_stress_test}
\footnotesize
\setlength{\tabcolsep}{2.5pt} % Slightly reduced to ensure fit in portrait
\begin{tabularx}{\textwidth}{c cccc cc ccc cc}
\toprule
& \multicolumn{3}{c}{\textbf{Objective Values}} & \multicolumn{2}{c}{\textbf{VMC Metrics}} & \multicolumn{3}{c}{\textbf{Probabilities}} & \multicolumn{2}{c}{\textbf{Arcs}} & \\
\cmidrule(lr){2-4} \cmidrule(lr){5-6} \cmidrule(lr){7-9} \cmidrule(lr){10-11}
\thead{Amp.\\Factor} & \thead{Corr.} & \thead{Ind.} & \thead{Ind/Corr} & \thead{Abs.} & \thead{\%} & \thead{Joint\\Corr} & \thead{Joint\\Ind} & \thead{Ratio} & \thead{Corr} & \thead{Ind} & \thead{Diff?}\\
\midrule
1.00  & 47417.02 & 48094.16 & 47417.02 & 0.00 & 0.00 & 0.30 & 0.30 & 1.00 & 13 & 13 & F \\
1.50  & 48596.67 & 49170.39 & 48596.67 & 0.00 & 0.00 & 0.39 & 0.37 & 1.05 & 13 & 13 & F \\
2.00  & 49504.09 & 49992.23 & 49504.09 & 0.00 & 0.00 & 0.46 & 0.43 & 1.06 & 13 & 13 & F \\
3.00  & 50808.51 & 51164.46 & 50808.51 & 0.00 & 0.00 & 0.56 & 0.53 & 1.07 & 13 & 13 & F \\
5.00  & 52350.09 & 52535.87 & 52350.09 & 0.00 & 0.00 & 0.68 & 0.64 & 1.06 & 13 & 13 & F \\
8.00  & 53543.58 & 53587.23 & 53543.58 & 0.00 & 0.00 & 0.77 & 0.74 & 1.04 & 13 & 13 & F \\
10.00 & 54016.67 & 54001.48 & 54016.67 & 0.00 & 0.00 & 0.81 & 0.78 & 1.04 & 13 & 13 & F \\
15.00 & 54721.76 & 54616.22 & 54721.76 & 0.00 & 0.00 & 0.87 & 0.84 & 1.03 & 13 & 13 & F \\
\bottomrule
\end{tabularx}
\end{table}
 
Table~\ref{tab:tab13_stress_test} and Figure~\ref{fig:vmc-stress}
report the stress test. This persistent zero VMC reveals a structural
property of the calibrated network: the diversified route-and-inventory
portfolio that hedges marginal corridor disruptions automatically
provides protection against correlated multi-corridor events. The
design activates enough Cape alternatives, direct routes, and pipeline
bypasses to absorb the recourse impact of joint failures without
additional cost. We label this property \emph{structural joint-failure
resilience} and offer it as a candidate descriptor of the regime
occupied by the calibrated Indian network: VMC$\,=0$ on this network
because the marginal-risk-optimal design already activates routes that
bypass every corridor participating in a correlated disruption
scenario, so the recourse value is invariant to how the marginal
failure probabilities are coupled. This is a positive finding rather
than a limitation: it identifies a calibration regime in which the
simpler independent-marginal model suffices. Whether structural
joint-failure resilience holds on different network topologies, under
tighter capacity constraints, or under more binding shortage
penalties is an empirical question that we leave open.



The exact-zero VMC across both $\rho$-interpolation and $\gamma$-amplification stress tests is a substantive managerial finding rather than a calibration artifact. The Indian maritime network admits a first-stage design—crude pipelines via UAE ADCOP and Saudi Petroline, strategic LNG and LPG inventory at Dahej, Mumbai, and Vizag, and a Cape diversification leg—whose recourse value is invariant to how Hormuz and Bab/Suez correlate, because the activated route portfolio absorbs joint failure through the same hedging mechanism that absorbs single failures. The structural reading is that route-portfolio diversification dominates correlation calibration: once an adequately diversified portfolio is selected, the marginal value of resolving joint-failure dependence drops to zero. This finding has direct policy value. Calibrating maritime-disruption joint-probability tables from sparse historical co-occurrence data is notoriously hard; the result here implies that, in the regime occupied by India, planners can tolerate substantial uncertainty in joint-disruption probabilities without operational consequence, provided the route portfolio is chosen via a stochastic program that internalizes the marginal disruption distribution.


\subsection{Value of Decision-Dependent Probabilities}
\label{sec:vep}
 
The previous experiment showed that the calibrated base-case design is
structurally resilient to joint-failure correlation. We now isolate the
role of the decision-dependent probability mechanism. In the proposed
model, scenario probabilities depend on the first-stage arc-activation
vector through
\[
p_s(y)=\bar p_s+\sum_{a\in A}\delta_{sa}y_a.
\]
This captures endogenous exposure: a design that relies more heavily on
vulnerable corridors should be evaluated under a probability
distribution that reflects that exposure.
 
To isolate this effect, we compare a DDU-aware model with a
fixed-probability counterfactual. Both models use the same correlated
scenario library and are solved in a risk-neutral configuration, with
CVaR disabled. The DDU-aware model optimizes using the exposure-adjusted
probabilities \(p_s(y)\), while the no-DDU model optimizes using the
baseline probabilities \(\bar p_s\). We then evaluate the no-DDU
first-stage design under the DDU-aware model.
 
Let \(Z_{\mathrm{DDU}}\) denote the optimal objective under the
DDU-aware model, and let \(Z_{\mathrm{NoDDU}\rightarrow\mathrm{DDU}}\)
denote the objective obtained by evaluating the fixed-probability
design under the DDU-aware probability model. We define the value of
endogenous probabilities as
\[
\mathrm{VEP}
=
Z_{\mathrm{NoDDU}\rightarrow\mathrm{DDU}}
-
Z_{\mathrm{DDU}}.
\]
A positive VEP indicates that ignoring endogenous exposure leads to a
first-stage design that performs worse when evaluated under
decision-dependent scenario probabilities.
 
\begin{table}
\caption{Value of endogenous probabilities.}
\label{tab:tab04_vdu}
\begin{tabular}{rrrrrr}
\toprule
DDU\_objective & NoDDU\_objective & NoDDU\_design\_under\_DDU & VEP & VEP\_pct & ordering\_holds \\
\midrule
47225.67 & 47417.02 & 47663.03 & 437.37 & 0.93 & True \\
\bottomrule
\end{tabular}
\end{table}

 
Table~\ref{tab:tab04_vdu} reports the value of decision-dependent
probabilities at the calibrated DDU magnitude (\(k=1\)). The DDU-aware
model has objective value \(47{,}225.67\), whereas the fixed-probability
model has objective value \(47{,}417.02\) under its own
baseline-probability evaluation. When the fixed-probability design is
evaluated under the DDU-aware model, its objective increases to
\(47{,}663.03\). The resulting VEP is \(437.37\), or \(0.93\%\) of the
DDU-aware objective.
 
\paragraph{Sensitivity of VEP to DDU magnitude}
 
The calibrated \(\delta_{sa}\) coefficients are bounded by 0.005 per arc
(see Section~3), and the resulting probability shifts are small. To
understand how VEP responds to stronger exposure sensitivity, we
scale the entire \(\delta\) matrix by a multiplier
\(k\in\{0,0.5,1,2,3,5,k_{\max}\}\), where \(k_{\max}=6.7\) is the
largest multiplier consistent with the probability-validity constraint
\(p_s(y)\ge 0\) for all binary~\(y\).
 
\begin{table}[htbp] % Changed from sidewaystable
\centering
\caption{Parametric VEP sweep over DDU magnitude multiplier $k$.}
\label{tab:tab12_vep_sweep}
\scriptsize % Changed from footnotesize to fit 16 columns
\setlength{\tabcolsep}{1.5pt} % Reduced spacing further
\begin{tabularx}{\textwidth}{c ccc cc cc c cc cc c cc}
\toprule
& \multicolumn{3}{c}{\textbf{Objectives}} & \multicolumn{2}{c}{\textbf{VEP Metrics}} & \multicolumn{3}{c}{\textbf{Design (Arcs)}} & \multicolumn{3}{c}{\textbf{Inventory}} & \multicolumn{2}{c}{\textbf{Prob. Shift}} & \multicolumn{2}{c}{\textbf{Recourse}} \\
\cmidrule(lr){2-4} \cmidrule(lr){5-6} \cmidrule(lr){7-9} \cmidrule(lr){10-12} \cmidrule(lr){13-14} \cmidrule(lr){15-16}
$k$ & \thead{DDU\\Obj.} & \thead{NoDDU\\Obj.} & \thead{NoDDU\\/DDU} & \thead{Abs.} & \thead{\%} & \thead{DDU} & \thead{NoDDU} & \thead{Diff} & \thead{DDU} & \thead{NoDDU} & \thead{Diff} & \thead{Max} & \thead{Mean} & \thead{DDU} & \thead{NoDDU} \\
\midrule
0.00 & 47417 & 47417 & 47417 & 0.00 & 0.00 & 13 & 13 & 0 & 4250 & 4250 & 0 & 0.00 & 0.00 & 45881 & 45881 \\
0.50 & 47340 & 47417 & 47540 & 199.75 & 0.42 & 21 & 13 & 8 & 4250 & 4250 & 0 & 0.00 & 0.00 & 45788 & 45881 \\
1.00 & 47225 & 47417 & 47663 & 437.37 & 0.93 & 21 & 13 & 8 & 4250 & 4250 & 0 & 0.00 & 0.00 & 45673 & 45881 \\
2.00 & 46953 & 47417 & 47909 & 956.03 & 2.04 & 20 & 13 & 7 & 4250 & 4250 & 0 & 0.01 & 0.00 & 45403 & 45881 \\
3.00 & 46643 & 47417 & 48155 & 1512.04 & 3.24 & 19 & 13 & 6 & 4250 & 4250 & 0 & 0.02 & 0.01 & 45095 & 45881 \\
5.00 & 45946 & 47417 & 48647 & 2700.96 & 5.88 & 19 & 13 & 6 & 4250 & 4250 & 0 & 0.04 & 0.01 & 44398 & 45881 \\
6.70 & 45353 & 47417 & 49065 & 3711.54 & 8.18 & 19 & 13 & 6 & 4250 & 4250 & 0 & 0.05 & 0.01 & 43805 & 45881 \\
\bottomrule
\end{tabularx}
\end{table}
 
Table~\ref{tab:tab12_vep_sweep} reports the full sweep. VEP grows
monotonically from zero at \(k=0\) (DDU disabled) to \(3{,}711.54\) at
\(k=6.7\), corresponding to \(8.18\%\) of the DDU-aware objective. The
monotone increase confirms that VEP is not inherently small—it is
proportional to the planner's exposure sensitivity as encoded in the
\(\delta\) coefficients.
 
\begin{figure}[t]
    \centering
    \includegraphics[width=0.95\textwidth]{plots/fig12_vep_vs_k.pdf}
    \caption{Value of endogenous probabilities as a function of
    the DDU magnitude multiplier~\(k\). Left: VEP grows monotonically from
    zero to 8.18\% of objective. Right: DDU-aware design activates
    significantly more arcs than the no-DDU counterfactual, with
    diversification declining at high~\(k\); total pre-positioned inventory remains constant at 4{,}250 units across both models.}
    \label{fig:vep-sweep}
\end{figure}
 
Figure~\ref{fig:vep-sweep} reveals the mechanism. At \(k=0\), the
DDU-aware and no-DDU designs are identical: both activate 13 arcs. As
soon as \(k>0\), the DDU-aware planner activates 21 arcs—a
substantial diversification response—while the no-DDU planner
remains at 13 arcs across all~\(k\). Interestingly, the DDU-aware
design then \emph{reduces} its arc count from 21 to 19 as \(k\)
increases beyond~2. This non-monotone response reflects a shift in the
dominant hedging strategy: at moderate exposure sensitivity, the
planner diversifies across corridors to dilute per-corridor exposure;
at high sensitivity, it becomes cheaper to avoid vulnerable corridors
entirely rather than diversify across them. The DDU mechanism operates entirely through the route-selection channel: total pre-positioned inventory remains at 4{,}250 units for both models across all~\(k\) (cf.\ Table~\ref{tab:tab12_vep_sweep}). The exposure-adjusted probability model therefore changes \emph{which routes the planner contracts}, not \emph{how much the planner stores}.
 
The 0.93\% VEP at the calibrated $k=1$ is small in absolute terms but is bounded above only by the conservative magnitude of the analyst-set $\delta_{sa}$ coefficients (capped at $0.005$ per arc). The VEP-vs-$k$ sweep is approximately linear: each unit of $k$ adds roughly $1.2\%$ of objective in VEP, and at the probability-validity boundary $k_{\max}=6.7$ the VEP reaches $8.18\%$ of the DDU-aware objective. The interpretation is that, in this calibration, the measured value of the DDU mechanism is primarily governed by the exposure-sensitivity scale of $\delta$. The 0.93\% number at $k=1$ reflects the analyst's prior; an analyst who treats $\delta$ as known would conclude DDU is unimportant for India, whereas an analyst willing to assume an empirically grounded $\delta$ would reach the opposite conclusion. This separation makes $\delta$ the central calibration object of the formulation, although the precise VEP for any given $\delta$ also depends on network topology, shortage-penalty levels, route substitutability, and first-stage activation costs. Two consequences follow. First, sensitivity analysis around $k$ (Figure~\ref{fig:vep-sweep}) is the appropriate object for the planner's decision, not the point estimate at $k=1$. Second, the natural next step is to learn $\delta$ from observable geopolitical and maritime-intelligence data streams (event-coded conflict indicators, war-risk insurance premiums, AIS-derived transit volumes) rather than to set it via expert judgment; a data-driven $\delta$ calibration framework is left to future work, as discussed in Section~\ref{sec:conclusion}.


\subsection{Risk Aversion and the Mean--CVaR Frontier}
\label{sec:cvar-frontier}
 
Having isolated the value of endogenous probabilities, we next examine
the third modeling layer: risk aversion. The proposed objective combines
expected recourse cost and CVaR through the parameter \(\lambda\):
\[
(1-\lambda)\mathbb{E}[Q_s] + \lambda \mathrm{CVaR}_{\alpha}.
\]



When \(\lambda=0\), the model reduces to risk-neutral stochastic
planning. As \(\lambda\) increases, the planner places more weight on
tail-disruption losses. We first sweep $\lambda$ at fixed $\alpha=0.95$.

\begin{table}[ht]
\centering
\caption{Mean-CVaR frontier ($\lambda$ sweep).}
\label{tab:tab05_frontier}
\small
\setlength{\tabcolsep}{3pt} 
\begin{tabularx}{\textwidth}{c c CCCC CCC}
\toprule
& & \multicolumn{4}{c}{\textbf{Optimization Metrics}} & \multicolumn{3}{c}{\textbf{Design \& Performance}} \\
\cmidrule(lr){3-6} \cmidrule(lr){7-9}
$\lambda$ & $\alpha$ & \thead{Total\\Obj.} & \thead{Fixed\\Cost} & \thead{Exp.\\Recourse} & \thead{CVaR\\(Risk)} & \thead{Arcs\\(\#)} & \thead{Total\\Inv.} & \thead{Time\\(s)} \\
\midrule
0.00 & 0.95 & 47226 & 1552 & 45674 & 70212 & 21 & 4250 & 0.05 \\
0.10 & 0.95 & 49679 & 1552 & 45674 & 70212 & 21 & 4250 & 0.05 \\
0.20 & 0.95 & 52133 & 1552 & 45674 & 70212 & 21 & 4250 & 0.05 \\
0.30 & 0.95 & 54587 & 1552 & 45674 & 70212 & 21 & 4250 & 0.04 \\
0.40 & 0.95 & 57041 & 1552 & 45674 & 70212 & 21 & 4250 & 0.04 \\
0.50 & 0.95 & 59495 & 1552 & 45674 & 70212 & 21 & 4250 & 0.04 \\
0.60 & 0.95 & 61949 & 1552 & 45674 & 70212 & 21 & 4250 & 0.05 \\
0.70 & 0.95 & 64402 & 1552 & 45674 & 70212 & 21 & 4250 & 0.04 \\
0.80 & 0.95 & 66856 & 1552 & 45674 & 70212 & 21 & 4250 & 0.05 \\
\midrule
0.90 & 0.95 & 69310 & 1550 & 45693 & 70212 & 20 & 4250 & 0.05 \\
1.00 & 0.95 & 71736 & 1524 & 58364 & 70212 & 7  & 4250 & 0.01 \\
\bottomrule
\end{tabularx}
\end{table}

Table~\ref{tab:tab05_frontier} reports a nearly flat CVaR response:
$\mathrm{CVaR}_{0.95}$ changes by only $0.125$ as $\lambda$ varies from
$0$ to $1$. We now show that this flatness is an artifact of the
interaction between the discrete nine-scenario tail and the high
confidence level $\alpha=0.95$, not an intrinsic property of the
formulation.
 
\paragraph{Sensitivity to the confidence level \(\alpha\).}
 
We fix \(\lambda=0.5\) and solve the full DDU-aware model for
\(\alpha\in\{0.50,0.60,0.70,0.80,0.85,0.90,0.95,0.99\}\).
 
\begin{table}[ht]
\centering
\caption{CVaR $\alpha$ sweep at $\lambda = 0.5$.}
\label{tab:tab14_alpha_sweep}
\small
\setlength{\tabcolsep}{3pt} % Tighten spacing to fit 11 columns
\begin{tabularx}{\textwidth}{c c CCC CC CCC}
\toprule
& & \multicolumn{3}{c}{\textbf{Costs}} & \multicolumn{2}{c}{\textbf{Risk Metrics}} & \multicolumn{3}{c}{\textbf{Design}} \\
\cmidrule(lr){3-5} \cmidrule(lr){6-7} \cmidrule(lr){8-10}
$\alpha$ & $\lambda$ & \thead{Total\\Obj.} & \thead{Fixed\\Cost} & \thead{Exp.\\Rec.} & \thead{CVaR} & \thead{VaR\\Thresh.} & \thead{Arcs\\(\#)} & \thead{Inv.\\Total} & \thead{Tail\\Scen.} \\
\midrule
0.50 & 0.50 & 54318 & 1550 & 45693 & 59842 & 47593 & 20 & 4250 & 4 \\
0.60 & 0.50 & 55849 & 1550 & 45693 & 62905 & 47593 & 20 & 4250 & 4 \\
0.70 & 0.50 & 56747 & 1550 & 45693 & 64701 & 58171 & 20 & 4250 & 3 \\
0.80 & 0.50 & 57827 & 1552 & 45674 & 66876 & 63187 & 21 & 4250 & 2 \\
0.85 & 0.50 & 58442 & 1552 & 45674 & 68106 & 63187 & 21 & 4250 & 2 \\
0.90 & 0.50 & 59009 & 1552 & 45674 & 69241 & 65953 & 21 & 4250 & 1 \\
0.95 & 0.50 & 59495 & 1552 & 45674 & 70212 & 70212 & 21 & 4250 & 0 \\
0.99 & 0.50 & 59495 & 1552 & 45674 & 70212 & 70212 & 21 & 4250 & 0 \\
\bottomrule
\end{tabularx}
\end{table}
 
Table~\ref{tab:tab14_alpha_sweep} reveals substantial CVaR variation.
At \(\alpha=0.50\), CVaR equals \(59{,}842\), reflecting the average cost
of the worst half of scenarios. At \(\alpha=0.95\), CVaR reaches
\(70{,}212\), the cost of the single worst scenario. The range of
\(10{,}370\) units (17.3\% of the \(\alpha=0.50\) value) shows that the
CVaR component is highly sensitive to \(\alpha\). The column
``n\_tail\_scenarios'' in the table explains the mechanism: at
\(\alpha=0.50\), four scenarios exceed the VaR threshold; at
\(\alpha=0.90\), only one scenario remains in the tail; at
\(\alpha\ge 0.95\), the tail degenerates to a single scenario whose
cost cannot be reduced by any first-stage design. The near-invariance of
CVaR to \(\lambda\) at \(\alpha=0.95\) (Table~\ref{tab:tab05_frontier})
is therefore a consequence of the discrete tail containing a single
immovable point, not a sign that risk aversion is irrelevant.
 
\begin{figure}[t]
    \centering
    \includegraphics[width=0.95\textwidth]{plots/fig14_alpha_sweep.pdf}
    \caption{Response to \(\alpha\). Left: mean--CVaR objective.
    Center: CVaR and expected recourse components. Right: design
    response showing a phase transition at \(\alpha\approx 0.75\),
    where the planner activates an additional arc.}
    \label{fig:alpha-sweep}
\end{figure}
 
Figure~\ref{fig:alpha-sweep} shows a design phase transition at
\(\alpha\approx 0.75\): the planner switches from 20 to 21 activated
arcs. Below this threshold, the tail contains enough scenarios that
the planner can manage tail risk with a leaner route portfolio. Above
it, the single remaining tail scenario (dual disruption) is severe
enough to justify activating an additional bypass arc. This threshold
effect has practical significance: it identifies the confidence level
at which the marginal cost of an additional hedging arc is justified
by the reduction in tail-risk cost.
 
\paragraph{Joint sensitivity to \((\lambda,\alpha)\)}
 
\begin{figure}[t]
    \centering
    \includegraphics[width=0.75\textwidth]{plots/fig14_lambda_alpha_heatmap.pdf}
    \caption{Mean--CVaR objective on a \(6\times 6\)
    \((\lambda,\alpha)\) grid. The \(\lambda=0\) row is constant
    (risk-neutral, \(\alpha\)-independent). Columns \(\alpha\ge 0.95\)
    saturate because the discrete tail contains a single scenario.}
    \label{fig:lambda-alpha}
\end{figure}
 
Figure~\ref{fig:lambda-alpha} reports the objective on a
\(6\times 6\) grid over \(\lambda\in\{0,0.2,0.4,0.6,0.8,1\}\) and
\(\alpha\in\{0.50,0.70,0.80,0.90,0.95,0.99\}\). Three features are
visible. First, the \(\lambda=0\) row is perfectly flat at \(47{,}226\):
the risk-neutral planner is by construction insensitive to~\(\alpha\).
Second, the rightmost two columns (\(\alpha=0.95\) and \(\alpha=0.99\))
produce identical values at every \(\lambda\), confirming the
single-scenario tail saturation. Third, the objective range across the
grid spans \(24{,}510\) units—a 52\% variation from the minimum
(\(47{,}226\) at \(\lambda=0\)) to the maximum (\(71{,}736\) at
\(\lambda=1,\alpha\ge 0.95\)). This shows that the mean--CVaR
formulation generates a rich set of Pareto-efficient designs when both
parameters are varied, even though the \(\lambda\)-only frontier at
fixed \(\alpha=0.95\) appears degenerate.
 
Overall, the \(\alpha\) sweep resolves the apparent flatness of the
base-case CVaR frontier. The formulation is not degenerate: the flat
frontier at \(\alpha=0.95\) reflects the finite-scenario tail structure,
and meaningful design variation emerges when \(\alpha\) is varied. For
policy applications with the nine-scenario library, \(\alpha\) values
in the range \([0.70,0.90]\) produce the most informative risk--cost
tradeoffs, because the tail at these levels contains two to three
scenarios rather than one.




\subsection{Scalability and Benders Decomposition}
\label{sec:benders-scaling}
 
The preceding experiments isolate the value of the model's uncertainty
and risk components. We now turn to computational scalability. The
calibrated 16-node, 28-arc topology is small enough that the
extensive-form MILP is the natural solver of choice, and the role of
the Benders decomposition introduced in
\Cref{sec:algorithm} is therefore not to outperform the extensive form
but to provide a structurally consistent decomposition that exactly
recovers the extensive-form solution and creates a path toward larger
scenario or topology regimes that arise under sample-average
approximation, stress testing, or data-driven scenario generation.
Because the extensive form replicates second-stage flow, inventory, and
shortage variables for each scenario, its size grows linearly with the
number of scenarios; the Benders master, by contrast, keeps the
first-stage design variables and separates the recourse evaluation
across scenarios.
 
We evaluate the decomposition in two stages: first, validation and
scenario-scaling on the calibrated network with expanded scenario
libraries; second, convergence analysis and a diagnostic on the
corridor-based valid inequalities.
 
\paragraph{Scenario scaling on the calibrated network}

We compare the monolithic extensive form and Benders decomposition on
expanded scenario libraries with \(|S|\in\{9,27,81,243,729\}\). To
isolate algorithmic scaling from the additional nonlinearities introduced
by risk aversion and decision-dependent probabilities, this experiment
uses the risk-neutral fixed-probability configuration. The extensive form
provides an exact benchmark for validating the decomposition.

\begin{table}[htbp] % Changed from sidewaystable
\centering
\caption{Scenario scaling: EF vs. Benders (BD) on calibrated network. For $(|S|\in \{27,81,243\})$, \texttt{make\_scaled\_instance} replicates the calibrated 9-scenario library with renormalized probabilities; the probability-weighted objective is invariant under this replication, which explains the identical objective values across these three rows. For $(|S|=729)$, the scaling procedure introduces additional scenarios beyond pure replication, producing the slightly different objective $27{,}227.32$.}
\label{tab:tab15_scenario_scaling}
\scriptsize 
\setlength{\tabcolsep}{1.5pt} % Minimal padding to fit 14 columns
\begin{tabularx}{\textwidth}{r cccc ccccccc cc cc}
\toprule
& \multicolumn{4}{c}{\textbf{Extensive Form (EF)}} & \multicolumn{7}{c}{\textbf{Benders Decomposition (BD)}} & \multicolumn{2}{c}{\textbf{Comparison}} \\
\cmidrule(lr){2-5} \cmidrule(lr){6-12} \cmidrule(lr){13-14}
\thead{Scen.\\($N$)} & \thead{Obj.} & \thead{Time\\(s)} & \thead{Vars} & \thead{Cons} & \thead{Obj.} & \thead{Time\\(s)} & \thead{Iter.} & \thead{Cuts} & \thead{Gap\\\%} & \thead{Stat.} & \thead{Abs.\\Gap} & \thead{Faster?} & \thead{Speed\\up} \\
\midrule
9    & 47417.02 & 0.01 & 2261   & 1620   & 47417.02 & 0.10 & 11 & 67   & 0.00 & OPT & 0.0e+00  & F & 0.09 \\
27   & 25948.35 & 0.02 & 6599   & 4860   & 25948.35 & 0.21 & 11 & 234  & 0.00 & OPT & 7.3e-12  & F & 0.09 \\
81   & 25948.35 & 0.05 & 19613  & 14580  & 25948.35 & 0.58 & 10 & 648  & 0.00 & OPT & 1.1e-11  & F & 0.08 \\
243  & 25948.35 & 0.15 & 58655  & 43740  & 25948.35 & 1.71 & 10 & 1944 & 0.00 & OPT & 0.0e+00  & F & 0.09 \\
729  & 27227.32 & 0.48 & 175781 & 131220 & 27227.32 & 5.36 & 10 & 5670 & 0.00 & OPT & 6.9e-11  & F & 0.09 \\
\bottomrule
\end{tabularx}
\end{table}

Table~\ref{tab:tab15_scenario_scaling} reports the experiment. Across the
entire range, Benders converges to the EF optimum at numerical precision
(final absolute objective gap $\le 7\times 10^{-11}$ in every row),
validating the decomposition. Two algorithmic regularities are visible.
First, the iteration count is essentially constant: Benders terminates in
10 or 11 iterations in every row, independent of $|S|$. This is the
empirical analogue of \Cref{prop:benders_finite}: convergence in the
iteration metric is governed by the dual structure of the scenario
subproblems, not by their cardinality. Second, the cut count grows
linearly with $|S|$: the sequence $67/234/648/1944/5670$ is closely
approximated by $8|S|$, reflecting the per-scenario optimality cuts
added at each iteration.

EF outperforms Benders by an essentially constant factor (approximately
$11\times$) across the entire scenario range tested. This is the expected
scaling regime for small geometric instances: with only 16 nodes and 28
arcs, the EF MIP relaxation is sufficiently small that monolithic
solution dominates iterative cut generation, and adding scenarios on a
fixed topology does not change that ranking. The advantage of
decomposition emerges along the topological dimension rather than the
scenario dimension; preliminary tests on random S-CMNDP instances of up
to 80 nodes and 602 arcs (omitted for brevity) show the EF-to-Benders
runtime ratio narrowing toward $1\times$ as the EF model size approaches
$10^6$ variables.

\paragraph{Benders convergence and group-failure cut diagnostic}
We next examine convergence behavior and the effect of the corridor-based
group-failure cuts on the base stochastic instance. Without the group
cuts, the algorithm terminates in 11 iterations with final lower and
upper bounds both equal to \(47{,}417.02\), yielding a final optimality
gap of \(0\%\) and a total of 67 added cuts.

\begin{table}
\caption{Benders convergence trace.}
\label{tab:tab10_benders_convergence}
\begin{tabular}{rrrrrrr}
\toprule
iteration & lower\_bound & upper\_bound & gap\_abs & gap\_pct & cuts\_added & runtime\_sec \\
\midrule
1 & 0.00 & 256500.00 & 256500.00 & 100.00 & 9 & 0.02 \\
2 & 2785.51 & 54272.94 & 51487.43 & 94.87 & 9 & 0.03 \\
3 & 38916.94 & 54272.94 & 15356.00 & 28.29 & 9 & 0.03 \\
4 & 39049.55 & 54272.94 & 15223.40 & 28.05 & 9 & 0.04 \\
5 & 39669.82 & 54272.94 & 14603.12 & 26.91 & 9 & 0.05 \\
6 & 39671.82 & 47486.90 & 7815.08 & 16.46 & 7 & 0.06 \\
7 & 47284.90 & 47486.90 & 201.99 & 0.43 & 9 & 0.06 \\
8 & 47374.40 & 47486.90 & 112.50 & 0.24 & 2 & 0.07 \\
9 & 47376.40 & 47417.02 & 40.62 & 0.09 & 2 & 0.07 \\
10 & 47385.24 & 47417.02 & 31.78 & 0.07 & 2 & 0.08 \\
11 & 47417.02 & 47417.02 & 0.00 & 0.00 & 0 & 0.08 \\
\bottomrule
\end{tabular}
\end{table}


\begin{figure}[t]
    \centering
    \includegraphics[width=0.78\textwidth]{plots/fig10_benders_convergence.pdf}
    \caption{Benders convergence trace on the base instance. The lower
    bound is obtained from the master problem; the upper bound is
    obtained by evaluating incumbent first-stage designs through scenario
    subproblems.}
    \label{fig:benders-convergence}
\end{figure}

Figure~\ref{fig:benders-convergence} shows the bound trajectory; the
two bounds meet at iteration~11, confirming finite convergence on the
calibrated instance.

Adding the corridor-based group-failure cuts of \Cref{prop:group-cut}
preserves the optimum (47{,}417.02) and the iteration count (11), but
slightly increases wall-clock runtime (0.08\,s without cuts vs.\ 0.10\,s
with cuts). On this nine-scenario instance the additional cut-management
overhead dominates any strengthening benefit. This negative diagnostic
is informative rather than dismissive: corridor-based cuts encode valid
domain structure, but their computational value depends on instance
size, scenario richness, and the weakness of the master relaxation, none
of which favors strengthening on the calibrated nine-scenario instance.

\paragraph{Takeaway}
The decomposition experiments establish two points. First, Benders
matches the extensive-form objective to numerical precision on every
converged instance and terminates in 10--11 iterations independent of
$|S|$ up to $|S|=729$, validating the empirical content of
\Cref{prop:benders_finite}. Second, the unstrengthened decomposition
already closes the optimality gap quickly on the calibrated instance,
and the corridor-based group-failure cuts neither tighten the master
nor reduce wall-clock time. This second observation, together with the
linear cut-pool growth at $|S|=729$ (5{,}670 retained cuts versus 10
productive iterations), points to the operative algorithmic bottleneck
in larger settings: the binding question is not which valid cuts to
generate but which to retain. We return to this point in
\Cref{sec:conclusion}.



% Future work — promoted to a contributed section once implementation is complete.
% \section{Learning-Augmented Benders Decomposition}
% \label{sec:la-benders}

\section{Conclusion}
\label{sec:conclusion}

We have developed a two-stage stochastic multi-commodity flow model
for maritime energy supply chain resilience that integrates correlated
chokepoint disruptions, decision-dependent disruption probabilities,
and mean--CVaR risk aversion in a single tractable MILP. Three
structural results---an exact McCormick reformulation of the bilinear
DDU--CVaR objective, scenario-wise decomposability of the recourse
problem, and finite convergence of Benders decomposition---enable
solution by both extensive-form MILP and Benders decomposition with
correlation-exploiting cuts.

The computational study, calibrated to Indian maritime energy imports
under the 2026 Hormuz crisis, isolates the contribution of each
modeling layer. Stochastic modeling reduces the cost of an
implementable design by 14.8\% relative to the expected-value
approximation. On the calibrated Indian network and across both
$\rho$-interpolation and $\gamma$-amplification stress tests, the
value of modeling correlation is exactly zero, a substantive finding
rather than a null result: the diversified route-and-inventory
portfolio selected by the stochastic program is itself a sufficient
hedge against joint-corridor failure. We caution that this property
is a feature of the specific network topology, capacity profile, and
shortage-penalty structure of the calibrated case study; whether it
holds on tighter or qualitatively different networks is left as an
empirical question. The value of endogenous probabilities is 0.93\%
at the analyst-set exposure sensitivity and grows approximately
linearly to 8.18\% at the probability-validity boundary; in this
calibration, VEP is therefore primarily governed by the exposure
sensitivity scale of the DDU coefficients $\delta$ rather than by
the network or shortage-penalty parameters, identifying $\delta$ as the
central calibration object of the formulation. The mean--CVaR
frontier exhibits a design phase transition at $\alpha\approx 0.75$,
where the planner activates an additional bypass arc to hedge the
single remaining tail scenario.

At the commodity level, the optimal policy fully protects crude oil
and fertilizer through pipeline bypass and strategic reserves, while
LNG and LPG retain residual shortage exposure under severe
multi-corridor disruption; LPG is the most exposed commodity and the
natural target of any incremental resilience investment.

Three avenues are open for follow-on work. First, learning the
decision-dependent probability coefficients $\delta_{sa}$ from
observable geopolitical and maritime-intelligence data streams, rather
than setting them by analyst judgment, would convert the central
calibration object of the model into an empirically grounded one and
plausibly raise the realized VEP into the 5--8\% range identified by
the sensitivity analysis. Second, the cut-pool growth observed at
$|S|=729$ (5{,}670 cuts retained while only 10 iterations are
productive) suggests that learned cut-management procedures—scoring
candidate cuts by predicted contribution to bound tightening and
pruning low-utility cuts on the fly—could substantially accelerate
Benders at the scenario sizes that arise under sample-average
approximation or data-driven scenario generation. Third, the
structural joint-failure resilience finding for the Indian network
raises a comparative question: which other national maritime energy
networks exhibit the same structural property, and which do not, and
what topological or capacity features distinguish the two regimes.


\bibliographystyle{elsarticle-harv}
\bibliography{reference}

\appendix

%================================================================
\section{Institutional Mapping for the Indian Calibration}
\label{app:institutional}
%================================================================

This appendix documents the mapping between model variables and the
Indian institutional authorities used to calibrate the case study in
\Cref{sec:computational}. The mapping is intended for transparency of
the calibration; the formulation in \Cref{sec:model} is not specific to
any single institutional setting.

\paragraph{Strategic inventory ($W_{ic}$).}
First-stage inventory decisions correspond to the pre-positioning of
crude oil in underground rock caverns operated by the Indian Strategic
Petroleum Reserves Limited (ISPRL), which maintains 5.33~million metric
tonnes (MMT) of strategic crude across three coastal sites,
Visakhapatnam (1.33~MMT), Mangalore (1.5~MMT), and Padur (2.5~MMT),
providing approximately 9.5 days of national consumption cover
\citep{eia2016spr, pib2021spr}. For LNG, Petronet LNG Ltd.\ operates
India's largest regasification terminal at Dahej (17.5~MMTPA capacity)
with buffer storage; the Kochi terminal provides additional capacity on
the southwest coast. LPG and fertilizer storage is limited to
commercial tankage at port terminals.

\paragraph{Arc activation ($y_a$).}
Binary arc-activation decisions represent three classes of sovereign or
quasi-sovereign action: (i)~long-term government-to-government supply
contracts, exemplified by the Petronet--QatarEnergy LNG SPA (7.5~MMTPA
through 2048) \citep{petronet2024qatar}; (ii)~pipeline infrastructure
sanctioning, corresponding to the UAE Habshan--Fujairah ADCOP pipeline
(1.5~mb/d) and the Saudi East--West Petroline (7~mb/d, expanded 2025)
\citep{iea2026hormuz, cnbc2026pipelines}; and (iii)~chartering and
flag-state regulatory authority exercised by the Directorate General of
Shipping \citep{indiastrategic2026hormuz}, with the corridor-dependence
cap $\bar\Phi_\ell$ representing the regulatory ceiling on corridor
concentration.

\paragraph{Risk aversion ($\lambda$, $\alpha$) and DDU coefficients ($\delta_{sa}$).}
The CVaR weight and confidence level encode the planner's risk
tolerance, set in practice at the Cabinet Committee on Security level;
the 2026 Hormuz crisis activated multi-ministerial response including
LPG rationing and Essential Commodities Act enforcement
\citep{tribune2026hormuz}. The DDU shift parameters $\delta_{sa}$ are
exposure-adjusted planning probabilities reflecting forward-looking
intelligence assessments produced by the MEA, R\&AW, and the Joint
Intelligence Committee, rather than physical causal effects of arc
activation.

\paragraph{Generalizability.}
With minor relabeling, the same mapping applies to the United States
(SPR managed by the Department of Energy, with drawdown authority
resting with the President), China (NDRC coordinating SINOPEC, CNPC,
and CNOOC under a national energy security mandate), and the European
Union (mandatory gas storage filling under
Regulation~(EU)~2022/1032) \citep{eu2022gasstorage}.

\newpage
%================================================================
\section{Calibration Protocol}
\label{app:calibration}
%================================================================

This appendix documents the source and construction rule for every
numerical input to the calibrated case study. All values were retrieved
or last verified on 2026-04-18. The model carries dimensionless
normalised units (1 crude unit $\approx$ 1 kb/d; 1 LNG/LPG/fertilizer
unit $\approx$ 12 kt/yr) so that the four commodities share a single
linear program; the convention is chosen so that Indian aggregate
imports (4.5\,mb/d crude, 17.5\,Mt/yr LNG, 23.3\,Mt/yr LPG, 9.2\,Mt/yr
fertilizer feedstock) map to comparable per-arc capacities. Costs are
stored in abstract relative units, anchored so that the cheapest
Gulf$\rightarrow$India route (Bandar Imam $\rightarrow$ Jamnagar, 3
days) has cost $4.0$. \Cref{tab:calibration} summarises the
calibration; the full source listing is available in the project
repository.

% \begin{table}[H]
% \centering\small
% \caption{Calibration protocol for the Indian case study. Anchors are
%   national-aggregate quantities; calibrated values are the per-arc /
%   per-node / per-scenario inputs in the project CSV files.}
% \label{tab:calibration}
% \begin{tabularx}{\textwidth}{@{}p{2.6cm} p{3.3cm} p{3.0cm} X@{}}
% \toprule
% Parameter & Anchor & Source & Construction rule \\
% \midrule
% National crude import & 5.5 mb/d total; 4.5 mb/d via Hormuz (2024) &
%   EIA, IEA \emph{Strait of Hormuz} (2026) &
%   Apportioned across Jamnagar / Mumbai / Vizag in proportion to refinery
%   capacity. \\
% \addlinespace
% National LNG import & 17.5 Mt/yr; 69\% Hormuz-linked (2025) &
%   PPAC, Petronet LNG \citep{petronet2024qatar} &
%   Allocated 69\% Dahej, 31\% Kochi (terminal-capacity weighted). \\
% \addlinespace
% National LPG import & 23.3 Mt/yr; 97\% Middle East (2024) &
%   Drewry LPG outlook &
%   Mumbai, Vizag in proportion to coastal LPG terminal capacity. \\
% \addlinespace
% Fertilizer import & 7\,Mt urea + 2.2\,Mt ammonia (2023--24) &
%   PIB, FAI Statistical Database &
%   Paradip and Mumbai in proportion to port silo capacity. \\
% \addlinespace
% ISPRL crude storage $L_{ic}$ & 5.33 MMT (Phase I) &
%   PIB \citep{pib2021spr}, EIA \citep{eia2016spr} &
%   Vizag 1.33, Mangalore 1.5, Padur 2.5 (assigned to nearest demand
%   node). LNG storage anchored to terminal tank volumes (2$\times$160{,}000\,m$^3$
%   at Dahej; 2$\times$155{,}000\,m$^3$ at Kochi). \\
% \addlinespace
% Hormuz capacity & 20.7 mb/d (2024) &
%   EIA \emph{Hormuz Factsheet} \citep{iea2026hormuz} &
%   Allocated to Hormuz-routed arcs by bilateral trade share
%   (UNCTAD 2024). \\
% \addlinespace
% Bypass-pipeline capacity & ADCOP 1.5 mb/d; Petroline 7 mb/d (2025) &
%   IEA, CNBC \citep{iea2026hormuz, cnbc2026pipelines} &
%   Allocated to the two pipeline arcs at nameplate. \\
% \addlinespace
% Suez/Bab-el-Mandeb capacity & 4.9 mb/d Suez; 4.1--4.2 mb/d Bab (1H 2024--25) &
%   EIA \emph{Red Sea decline 2024} &
%   Allocated to corresponding bab\_suez arcs at fleet share. \\
% \addlinespace
% Transport cost $c_{ac}^{s}$ & WS70 baseline freight 2023--24 ($\sim$\$1.6/bbl Ras Tanura$\rightarrow$Jamnagar) &
%   Lloyd's List, Argus &
%   Anchored at $c=4.0$ for the cheapest 3-day Gulf$\rightarrow$India route;
%   all other base costs scaled by relative voyage time. Scenario costs
%   $c_{ac}^{s}=c_{ac}\cdot\textsf{cost\_mult}_{\ell s}$. \\
% \addlinespace
% Holding cost $h_c$ & Cryogenic / pressurised storage premia &
%   Industry storage-cost surveys &
%   Crude 0.3, LNG 0.8, LPG 0.5, fertilizer 0.2 (relative ranking by
%   storage-class technology). \\
% \addlinespace
% Shortage penalty $\pi_c$ & Value-of-lost-load + welfare-loss &
%   Outlook India 2026 (LPG); refinery feedstock value &
%   Crude 25, LNG 40, LPG 50, fertilizer 35. The LPG penalty is set above
%   LNG to reflect the household welfare and political-cost channel
%   documented during the 2026 Hormuz crisis. \\
% \addlinespace
% Activation cost $f_a$ & — & — &
%   Set at 2.0 for every candidate arc, calibrated so that total
%   activation budget does not dominate recourse cost in the base case. \\
% \addlinespace
% Corridor caps $\bar\Phi_\ell$ & DG Shipping diversification mandate &
%   \citep{indiastrategic2026hormuz} &
%   Hormuz 0.7, Bab/Suez 0.5, Cape 0.8, pipelines and direct routes 1.0
%   (caps tighter on chokepoint corridors). \\
% \addlinespace
% Baseline scenario probabilities $\bar p_s$ & IEA / EIA disruption-risk surveys; UNCTAD March-2026 brief &
%   \citep{iea2026hormuz}, EIA scenario discussion &
%   Nine scenarios: see \Cref{tab:tab01_scenario_library}. Weights
%   reflect single-corridor severity gradient and a calibrated
%   joint-disruption mass of 0.08 (scenario S5). \\
% \addlinespace
% Disruption multipliers & UKMTO 80\% Hormuz transit reduction; World Bank 75\% Suez decline; UNCTAD 42\% Cape distance increase &
%   UKMTO advisories Mar-2026; UNCTAD 2024 &
%   Capacity / cost / time multipliers per (corridor, scenario): see
%   \Cref{tab:tab01_scenario_library}. \\
% \addlinespace
% DDU coefficients $\delta_{sa}$ & — (analyst-set; sign anchored to corridor exposure) &
%   This paper &
%   $|\delta_{sa}|\le 0.005$, with $\sum_s \delta_{sa} = 0$ enforced
%   at construction; sign of $\delta_{sa}$ is positive for activations
%   that increase exposure to scenario $s$ and negative for activations
%   that reduce it. The magnitude is the smallest value satisfying
%   $p_s(y)\ge 0$ for every binary $y$; sensitivity to a uniform scaling
%   multiplier $k$ is reported in Section~\ref{sec:vep}. \\
% \addlinespace
% Risk parameters $(\lambda, \alpha)$ & — &
%   \citep{rockafellar2000} &
%   Base case $\lambda=0.5$, $\alpha=0.95$; sensitivity reported in
%   Section~\ref{sec:cvar-frontier}. \\
% \bottomrule
% \end{tabularx}
% \end{table}




% --- Revised Table Code ---
\setlength{\tabcolsep}{3pt} % Tightens horizontal space
\footnotesize               % Smaller font to fit more text
\begin{xltabular}{\textwidth}{@{} p{2.2cm} p{2.5cm} p{2.5cm} X @{}}
    \caption{Calibration protocol for the Indian case study.} \label{tab:calibration} \\
    \toprule
    \textbf{Parameter} & \textbf{Anchor} & \textbf{Source} & \textbf{Construction rule} \\
    \midrule
    \endfirsthead
    
    \multicolumn{4}{c}%
    {{\bfseries \tablename\ \thetable{} -- continued from previous page}} \\
    \toprule
    \textbf{Parameter} & \textbf{Anchor} & \textbf{Source} & \textbf{Construction rule} \\
    \midrule
    \endhead

    \bottomrule
    \endfoot

    National crude import & 5.5 mb/d total; 4.5 mb/d via Hormuz (2024) & EIA, IEA \emph{Hormuz} (2026) & Apportioned across Jamnagar / Mumbai / Vizag in proportion to refinery capacity. \\ \addlinespace
    
    National LNG import & 17.5 Mt/yr; 69\% Hormuz-linked (2025) & PPAC, Petronet \citep{petronet2024qatar} & Allocated 69\% Dahej, 31\% Kochi (terminal-capacity weighted). \\ \addlinespace
    
    National LPG import & 23.3 Mt/yr; 97\% Middle East (2024) & Drewry LPG outlook & Mumbai, Vizag in proportion to coastal LPG terminal capacity. \\ \addlinespace
    
    Fertilizer import & 7\,Mt urea + 2.2\,Mt ammonia (2023--24) & PIB, FAI Database & Paradip and Mumbai in proportion to port silo capacity. \\ \addlinespace
    
    ISPRL crude storage $L_{ic}$ & 5.33 MMT (Phase I) & PIB \citep{pib2021spr} & Vizag 1.33, Mangalore 1.5, Padur 2.5. LNG storage anchored to tank volumes (Dahej/Kochi). \\ \addlinespace
    
    Hormuz capacity & 20.7 mb/d (2024) & EIA \emph{Factsheet} \citep{iea2026hormuz} & Allocated to Hormuz-routed arcs by bilateral trade share. \\ \addlinespace
    
    Bypass-pipeline capacity & ADCOP 1.5 mb/d; Petroline 7 mb/d (2025) & IEA, CNBC \citep{cnbc2026pipelines} & Allocated to the two pipeline arcs at nameplate. \\ \addlinespace
    
    Suez/Bab capacity & 4.9 mb/d Suez; 4.1--4.2 mb/d Bab & EIA \emph{Red Sea 2024} & Allocated to corresponding bab\_suez arcs at fleet share. \\ \addlinespace
    
    Transport cost $c_{ac}^{s}$ & WS70 baseline freight ($\sim$\$1.6/bbl) & Lloyd's List, Argus & Costs scaled by relative voyage time. $c_{ac}^{s}=c_{ac}\cdot\textsf{cost\_mult}_{\ell s}$. \\ \addlinespace
    
    Holding cost $h_c$ & Storage premia & Industry surveys & Crude 0.3, LNG 0.8, LPG 0.5, fertilizer 0.2 (by tech class). \\ \addlinespace
    
    Shortage penalty $\pi_c$ & VOLL + welfare-loss & Outlook India 2026 & Crude 25, LNG 40, LPG 50, fertilizer 35. High LPG reflects household welfare risk. \\ \addlinespace
    
    Activation cost $f_a$ & — & — & Set at 2.0; calibrated so budget does not dominate recourse costs. \\ \addlinespace
    
    Corridor caps $\bar\Phi_\ell$ & Diversification mandate & \citep{indiastrategic2026hormuz} & Hormuz 0.7, Bab/Suez 0.5, Cape 0.8, Pipe 1.0. \\ \addlinespace
    
    Scenario probs $\bar p_s$ & IEA / EIA surveys; UNCTAD 2026 & \citep{iea2026hormuz} & Nine scenarios; joint-disruption mass of 0.08 (S5). \\ \addlinespace
    
    Disruption mult. & UKMTO / World Bank advisories & UKMTO, UNCTAD & Capacity / cost / time multipliers per corridor. \\ \addlinespace
    
    DDU coeffs $\delta_{sa}$ & Analyst-set & This paper & $|\delta_{sa}|\le 0.005$, $\sum \delta_{sa} = 0$. Sign based on corridor exposure. \\ \addlinespace
    
    Risk $(\lambda, \alpha)$ & — & \citep{rockafellar2000} & Base case $\lambda=0.5$, $\alpha=0.95$. \\
\end{xltabular}

\paragraph{Sensitivity coverage of the calibrated parameters.}
Of the parameters in \Cref{tab:calibration}, the paper reports
quantitative sensitivity for the DDU magnitude $k$
(Section~\ref{sec:vep}, full sweep to the probability-validity
boundary) and for both CVaR parameters $(\lambda, \alpha)$
(Section~\ref{sec:cvar-frontier}, including a full $6\times 6$ grid).
Sensitivity to shortage penalties $\pi_c$, corridor caps
$\bar\Phi_\ell$, and storage capacities $L_{ic}$ is not reported in
this paper; given the structural-resilience finding of
Section~\ref{sec:vmc} (VMC$\,=0$ across $\rho$- and $\gamma$-stress
tests), a full robustness sweep over these three parameters would
sharpen the boundary of the calibration regime in which the result
holds, but is computationally expensive and is left to follow-on work.

%================================================================
\section{Scenario Library and Base-Case Scenario Outcomes}
\label{app:basecase-additional}
%================================================================

This appendix reports the calibrated nine-scenario disruption library and
the scenario-level recourse outcomes under the base-case first-stage
design.

\paragraph{Scenario library.}
\Cref{tab:tab01_scenario_library} reports the disruption library used in
the case study. Each scenario specifies a joint disruption state over
the main maritime corridors and bypass options, together with the
multipliers determining how arc capacities, transport costs, and
transit-time-related surcharges change after the scenario is realized.
The probabilities are the baseline values $\bar p_s$ before any
decision-dependent adjustment.

\begin{sidewaystable}
\centering
\caption{Scenario library and disruption multipliers.}
\label{tab:tab01_scenario_library}
\small
\begin{tabularx}{\linewidth}{rlXXXXXXX}
\toprule
\textbf{ID} & \textbf{Scenario Name} & \textbf{Hormuz Cap.} & \textbf{Hormuz Cost} & \textbf{B-S Cap.} & \textbf{B-S Cost} & \textbf{Cape Cost} & \textbf{Pipe Cap.} & \textbf{Prob.} \\
\midrule
0 & baseline\_normal & 1.00 & 1.00 & 1.00 & 1.00 & 1.00 & 1.00 & 0.15 \\
1 & hormuz\_partial & 0.20 & 2.50 & 1.00 & 1.00 & 1.00 & 1.00 & 0.20 \\
2 & hormuz\_sev\_sel & 0.10 & 4.00 & 1.00 & 1.20 & 1.30 & 1.00 & 0.15 \\
3 & hormuz\_closure & 0.05 & 6.00 & 1.00 & 1.30 & 1.50 & 1.00 & 0.10 \\
4 & bab\_suez\_severe & 1.00 & 1.00 & 0.25 & 3.00 & 1.80 & 1.00 & 0.10 \\
5 & dual\_disruption & 0.15 & 3.50 & 0.20 & 3.50 & 2.00 & 1.00 & 0.08 \\
6 & insurance\_spike & 0.60 & 3.00 & 0.50 & 2.50 & 1.80 & 1.00 & 0.10 \\
7 & closure\_bypass & 0.05 & 6.00 & 0.80 & 1.50 & 1.60 & 1.50 & 0.07 \\
8 & delayed\_recovery & 0.40 & 2.00 & 0.60 & 1.80 & 1.40 & 1.00 & 0.05 \\
\bottomrule
\end{tabularx}
\end{sidewaystable}

\paragraph{Scenario-level outcomes under the base-case design.}
\Cref{tab:tab01_scenario_outcomes} reports the optimized scenario-level
recourse outcomes under the base-case first-stage design. For each
scenario, the table lists the baseline probability, the
decision-dependent probability evaluated at the optimized design
$p_s(y^*)$, the recourse cost $Q_s$, the shortage cost, and the number
of arcs carrying positive flow. These results are the numerical
counterpart to the scenario-cost figure in \Cref{sec:base-case} and
identify the Hormuz-related and dual-disruption scenarios as the
high-cost tail of the distribution.

\begin{sidewaystable}
\centering
\caption{Scenario-level base-case outcomes.}
\label{tab:tab01_scenario_outcomes}
\footnotesize % Slightly smaller to fit 11 columns
\begin{tabularx}{\linewidth}{rlXXXXXXXXX}
\toprule
\textbf{ID} & \textbf{Scenario} & \textbf{Prob.} & \textbf{DDU Prob.} & \textbf{Shift} & \textbf{Recourse Cost} & \textbf{CVaR Exc.} & \textbf{Short. Cost} & \textbf{Flow} & \textbf{Short. Qty} & \textbf{Arcs} \\
\midrule
0 & baseline & 0.15 & 0.15 & 0.00 & 19450.00 & 0.00 & 0.00 & 3750.00 & 0.00 & 9 \\
1 & hormuz\_part & 0.20 & 0.20 & 0.00 & 47592.75 & 0.00 & 17700.00 & 3370.00 & 380.00 & 8 \\
2 & hormuz\_sev & 0.15 & 0.15 & 0.00 & 58171.00 & 0.00 & 25100.00 & 3210.00 & 540.00 & 6 \\
3 & hormuz\_clos & 0.10 & 0.10 & -0.00 & 63187.00 & 0.00 & 27500.00 & 3150.00 & 600.00 & 5 \\
4 & bab\_suez & 0.10 & 0.10 & 0.00 & 19450.00 & 0.00 & 0.00 & 3750.00 & 0.00 & 9 \\
5 & dual\_disrupt & 0.08 & 0.08 & -0.00 & 70211.88 & 0.00 & 23900.00 & 3240.00 & 510.00 & 7 \\
6 & ins\_spike & 0.10 & 0.10 & 0.00 & 42709.25 & 0.00 & 2500.00 & 3700.00 & 50.00 & 9 \\
7 & clos\_bypass & 0.07 & 0.07 & 0.00 & 65953.00 & 0.00 & 27500.00 & 3150.00 & 600.00 & 5 \\
8 & del\_recovery & 0.05 & 0.05 & 0.00 & 39458.20 & 0.00 & 7900.00 & 3590.00 & 160.00 & 10 \\
\bottomrule
\end{tabularx}
\end{sidewaystable}

%================================================================
\section{Comparison Tables for the Value-of-Modeling Experiments}
\label{app:value-of-modeling-tables}
%================================================================

The figures in \Cref{sec:vss,sec:vmc,sec:vep} display the design
differences induced by the alternative formulations against the
baseline stochastic program. This appendix records the corresponding
tabulated comparisons used to construct those figures.

\paragraph{Stochastic program vs.\ expected-value design.}
\Cref{tab:tab02_design_comparison} compares the first-stage designs
obtained from the stochastic program (SP) and the expected-value (EV)
approximation. The differences in activated arcs and inventory levels
explain the cost gap measured by VSS in \Cref{sec:vss}.

\begin{table}[ht]
\centering
\caption{SP and EV first-stage design comparison.}
\label{tab:tab02_design_comparison}
\small
\begin{tabular}{lcccccccc}
\toprule
 & & & \multicolumn{5}{c}{\textbf{Pre-positioned Inventory}} \\
 \cmidrule(lr){4-8}
\textbf{Model} & \thead{True Obj.\\Value} & \thead{Arcs\\Active} & \textbf{Total} & \textbf{Crude} & \textbf{LNG} & \textbf{LPG} & \textbf{Fert.} \\
\midrule
SP & 47417.02 & 13 & 4250.00 & 2700.00 & 450.00 & 400.00 & 700.00 \\
EV & 54446.89 & 10 & 4250.00 & 2700.00 & 450.00 & 400.00 & 700.00 \\
\bottomrule
\end{tabular}
\end{table}

\paragraph{Correlated vs.\ independent-marginal probabilities.}
\Cref{tab:tab03_probability_comparison} compares the calibrated
correlated scenario probabilities with the independent-marginal
counterfactual. The independent counterfactual preserves marginal
corridor disruption probabilities but changes the probability assigned
to joint disruption states. Together with \Cref{tab:tab03_vmc} in the
main text, this confirms that the zero VMC reported in \Cref{sec:vmc}
arises not from probability identity but from design invariance under
the calibrated network topology.

\begin{table}[ht]
\centering
\caption{Correlated and independent-marginal scenario probabilities.}
\label{tab:tab03_probability_comparison}
\small
\begin{tabularx}{\textwidth}{rlCCCC}
\toprule
\textbf{ID} & \textbf{Scenario Name} & \thead{Correlated\\Prob.} & \thead{Indep.\\Prob.} & \thead{Prob.\\Diff.} & \thead{Abs.\\Diff.} \\
\midrule
0 & Baseline Normal & 0.15 & 0.15 & 0.00 & 0.00 \\
1 & Hormuz Partial & 0.20 & 0.15 & 0.05 & 0.05 \\
2 & Hormuz Severe/Sel. & 0.15 & 0.15 & 0.00 & 0.00 \\
3 & Hormuz Closure & 0.10 & 0.15 & -0.05 & 0.05 \\
4 & Bab Suez Severe & 0.10 & 0.10 & 0.00 & 0.00 \\
5 & Dual Disruption & 0.08 & 0.07 & 0.01 & 0.01 \\
6 & Insurance Spike & 0.10 & 0.07 & 0.03 & 0.03 \\
7 & Closure with Bypass & 0.07 & 0.07 & 0.00 & 0.00 \\
8 & Delayed Recovery & 0.05 & 0.07 & -0.02 & 0.02 \\
\bottomrule
\end{tabularx}
\end{table}

\paragraph{DDU-aware vs.\ fixed-probability design.}
\Cref{tab:tab04_design_comparison} compares the first-stage designs
optimized with and without decision-dependent probabilities. The
DDU-aware design activates substantially more arcs (21 versus 13) at
the same total inventory level, isolating route diversification as the
mechanism through which the DDU mechanism operates.

\begin{table}[htbp] % Changed from sidewaystable to standard table
\centering
\caption{DDU-aware and no-DDU first-stage design comparison.}
\label{tab:tab04_design_comparison}
\footnotesize 
\setlength{\tabcolsep}{2pt} % Adjusted for 15-column portrait fit
\begin{tabularx}{\textwidth}{l c c c ccccc cccccc}
\toprule
& & & & \multicolumn{5}{c}{\textbf{Pre-positioned Inventory}} & \multicolumn{6}{c}{\textbf{Activated Arcs by Corridor}} \\
\cmidrule(lr){5-9} \cmidrule(lr){10-15}
\textbf{Model} & \thead{DDU Obj.\\Value} & \thead{Total\\Arcs} & \thead{Total\\Inv.} & \textbf{Crude} & \textbf{LNG} & \textbf{LPG} & \textbf{Fert.} & \textbf{Util.} & \thead{Bab\\Suez} & \textbf{Cape} & \textbf{Dir.} & \textbf{Hor.} & \thead{Pipe\\KSA} & \thead{Pipe\\UAE} \\
\midrule
DDU   & 47225.67 & \textbf{21} & 4250 & 2700 & 450 & 400 & 700 & 1.00 & 0 & 6 & 6 & 7 & 1 & 1 \\
NoDDU & 47663.03 & 13 & 4250 & 2700 & 450 & 400 & 700 & 1.00 & 0 & 2 & 3 & 8 & 0 & 0 \\
\bottomrule
\end{tabularx}
\end{table}





\end{document}
